\documentclass[hidelinks,review,nohypdvips,onefignum,onetabnum]{siamart220329}
\makeatletter
\providecommand{\headerps@out}[1]{}
\makeatother

\usepackage{booktabs}
\usepackage{amssymb}
\usepackage{enumitem}
\usepackage{bm}
\usepackage{array}

% remark environment (not provided by siamart)
\newcounter{remark}[section]
\renewcommand{\theremark}{\thesection.\arabic{remark}}
\newenvironment{remark}[1][]{\refstepcounter{remark}\par\addvspace{\medskipamount}\noindent{\scshape Remark~\theremark\ifx&#1& \else\ (#1)\fi.}\ }{\par\addvspace{\medskipamount}}
\crefname{remark}{Remark}{Remarks}
\Crefname{remark}{Remark}{Remarks}

% Custom theorem environments (ntheorem-compatible, NOT \newtheorem)
\newcounter{assumption}[section]
\renewcommand{\theassumption}{\thesection.\arabic{assumption}}
\newenvironment{assumption}[1][]{\refstepcounter{assumption}\par\addvspace{\medskipamount}\noindent{\scshape Assumption~\theassumption\ifx&#1& \else\ (#1)\fi.}\itshape\ }{\par\addvspace{\medskipamount}}
\crefname{assumption}{Assumption}{Assumptions}
\Crefname{assumption}{Assumption}{Assumptions}

\newcounter{problem}[section]
\renewcommand{\theproblem}{\thesection.\arabic{problem}}
\newenvironment{problem}[1][]{\refstepcounter{problem}\par\addvspace{\medskipamount}\noindent{\scshape Problem~\theproblem\ifx&#1& \else\ (#1)\fi.}\itshape\ }{\par\addvspace{\medskipamount}}
\crefname{problem}{Problem}{Problems}
\Crefname{problem}{Problem}{Problems}

% Fix cleveref + hyperref + ntheorem interaction
\makeatletter
\def\cref@fix@refstepcounter@optarg[#1]#2{%
  \cref@old@refstepcounter{#2}%
  \cref@constructprefix{#2}{\cref@result}%
  \@ifundefined{cref@#1@alias}
    {\def\@tempa{#1}}
    {\def\@tempa{\csname cref@#1@alias\endcsname}}%
  \protected@edef\cref@currentlabel{%
    [\@tempa][\arabic{#2}][\cref@result]%
    \csname p@#2\endcsname\csname the#2\endcsname}}
\def\cref@fix@refstepcounter@noarg#1{%
  \cref@old@refstepcounter{#1}%
  \cref@constructprefix{#1}{\cref@result}%
  \@ifundefined{cref@#1@alias}
    {\def\@tempa{#1}}
    {\def\@tempa{\csname cref@#1@alias\endcsname}}%
  \protected@edef\cref@currentlabel{%
    [\@tempa][\arabic{#1}][\cref@result]%
    \csname p@#1\endcsname\csname the#1\endcsname}}
\AddToHook{begindocument/end}{%
  \let\refstepcounter@optarg\cref@fix@refstepcounter@optarg
  \let\refstepcounter@noarg\cref@fix@refstepcounter@noarg}
\makeatother

\newcommand{\N}{\mathbb{N}}
\newcommand{\R}{\mathbb{R}}
\newcommand{\cS}{\mathcal{S}}
\newcommand{\cU}{\mathcal{U}}
\newcommand{\cV}{\mathcal{V}}
\newcommand{\E}{\mathbb{E}}
\newcommand{\cK}{\mathcal{K}}
\newcommand{\cE}{\mathcal{E}}
\newcommand{\cF}{\mathcal{F}}
\newcommand{\cFN}{\mathcal{F}_N}
\newcommand{\cN}{\mathcal{N}}
\newcommand{\PhiT}{\Phi_{\theta}}
\newcommand{\diag}{\operatorname{diag}}
\newcommand{\softplus}{\operatorname{softplus}}
\newcommand{\norm}[1]{\left\lVert #1 \right\rVert}

\title{Deep Operator Learning of Maximum-Principle-Preserving Dissipative Semigroups for Time-Dependent PDEs}
\author{[Author Name]\thanks{[Affiliation]}}
\date{}

\begin{document}
\maketitle

\begin{abstract}
We formulate operator learning for autonomous evolution PDEs at the semigroup level: rather than learning only a one-step predictor, the learner is a family \(\Phi_{\theta}(\cdot,\tau)\) intended for repeated composition. We derive discrete stability estimates showing how trajectory-level approximation error and composition defect accumulate. We then construct two complementary finite-dimensional architectures. A bounded latent gradient flow provides interior-state invariance, exact composition of the continuous flow, and dissipation of an induced learned energy, conditionally on global existence. A hard affine boundary reconstruction and tangent autonomous Tanh generator define a globally well-posed forward semiflow that exactly preserves the tested discrete Dirichlet, Neumann, or Robin relation; explicit Runge--Kutta stages preserve the same relation. We derive a fourth-order nonuniform RK4 composition-defect estimate and a one-sided stability--generator-consistency bound that separates spatial, generator, initial-state, and time-integration errors. The PDE-level conclusion remains conditional on a PDE-specific theorem card: boundary consistency, method-of-lines convergence, physical invariants, and stability are not consequences of autonomy alone. Archived benchmark values are retained as pre-revision results, but they require rerunning under strict time alignment before supporting comparative accuracy claims.
\end{abstract}

\section{Introduction and semigroup formulation}\label{sec:introduction}
We study time-dependent PDEs whose solution operators form nonlinear semigroups and satisfy structural laws such as invariant-region preservation and Lyapunov dissipation. These laws are central both to the mathematical interpretation of the dynamics and to the stability of long-time repeated composition. Let \(X\) be a Banach space of states, for example \(L^{\infty}(\Omega)\), \(L^2(\Omega)\), or a Sobolev space on a spatial domain \(\Omega\subset \R^d\). Consider an autonomous evolution PDE of the form
\begin{equation}\label{eq:pde}
\partial_t u = \cF(u), \qquad t>0, \qquad u(0)=u_0\in X.
\end{equation}
Assume that \eqref{eq:pde} is well-posed \cite{CrandallLiggett1971} and generates a nonlinear solution semigroup
\begin{equation}\label{eq:semigroup-family}
\cS_t : X \to X, \qquad u(t)=\cS_tu_0, \qquad t\ge 0.
\end{equation}
Thus
\begin{equation}\label{eq:semigroup-law}
\cS_0 = \mathrm{Id}_X, \qquad \cS_{t+s}=\cS_t\circ \cS_s, \qquad t,s\ge 0.
\end{equation}
The autonomy assumption is essential to this one-parameter formulation.  If the generator depends explicitly on absolute time, the correct object is generally an evolution family \(U(t,s)\) satisfying a two-time cocycle law, as recorded in \eqref{eq:evolution-family-law}; it is not a semigroup on the unaugmented state.

Our focus is on problems for which the semigroup satisfies two additional structural properties.

\paragraph{Maximum principle / invariant region.}
There exists a closed admissible set \(\cK\subset X\) such that
\begin{equation}\label{eq:maximum-principle}
\cS_t(\cK)\subset \cK, \qquad t\ge 0.
\end{equation}
A typical example is
\begin{equation}\label{eq:box-invariant}
\cK = \{u\in X : m\le u(x)\le M \text{ a.e. in } \Omega\},
\end{equation}
which encodes a maximum principle, positivity, or boundedness property.

\paragraph{Dissipativity.}
There exists a Lyapunov functional
\begin{equation}\label{eq:energy-type}
\cE : X \to \R\cup\{+\infty\}
\end{equation}
for which
\begin{equation}\label{eq:dissipation}
\cE(\cS_tu) \le \cE(u), \qquad t\ge 0, \qquad u\in \cK.
\end{equation}
Depending on the PDE, \(\cE\) may be an energy, entropy, free energy, or another dissipated quantity.

Standard operator-learning formulations often emphasize short-time prediction. Our focus is different: we seek a learned evolution family intended for repeated composition in time. The approximation target is therefore the operator family \((\cS_t)_{t\ge 0}\) itself. This emphasis matters because a learned one-step map may fit short-time data and yet accumulate structural errors under repeated composition: small violations of invariance, dissipation, or semigroup consistency can be amplified when the model is used as a time-marching solver. The problem we address is to construct a learned family \((\PhiT(\cdot,\tau))_{\tau\ge 0}\) that approximates \((\cS_t)_{t\ge 0}\) while retaining \eqref{eq:maximum-principle} and \eqref{eq:dissipation}. In other words, the objective is not only local prediction accuracy but a learned evolution family that can be composed in time without destroying the PDE's structural laws.

\paragraph{Contributions.}
This paper makes five contributions.
First, it formulates structure-preserving operator learning at the nonlinear semigroup level: rather than learning an isolated one-step predictor, the target is an evolution family \(\Phi_{\theta}(\cdot,\tau)\) intended for repeated composition, with the PDE semigroup \((\cS_t)_{t\ge 0}\) as the reference object (Definition~\ref{def:main-learner}).
Second, it derives error-propagation estimates that decompose the two sources of error under repeated composition — trajectory-level approximation and semigroup defect — showing that both enter discrete time-stepping stability in complementary ways (Section~\ref{sec:repeated-composition}).
Third, it constructs a semidiscrete latent dissipative architecture (Sections~\ref{sec:latent-architecture}--\ref{sec:neural-param}) for which, assuming the latent ODE generates a global flow, invariance of the interior state set $(m,M)^N$, exact semigroup composition of the continuous flow, and dissipation of an induced learned energy hold by construction.
Fourth, \Cref{sec:boundary-semiflow} proves that the implemented hard affine reconstruction and tangent autonomous Tanh generator define a globally well-posed finite-dimensional semiflow preserving the tested discrete boundary relation.  It also proves explicit-RK boundary preservation, an RK4 partition-defect estimate, and a one-sided stability--consistency bound.
Fifth, \Cref{thm:transfer} gives a conditional exact-flow comparison with explicit projection and reconstruction maps, while \Cref{prop:numerical-transfer} adds numerical integration errors separately. \Cref{prop:universal-approx} supplies a proved error bound for architecture-compatible generators; density of the implemented restricted component classes remains open.
The layered character of the framework is deliberate: PDE, exact semidiscrete, exact latent, numerical, and empirical statements are not interchangeable.

\Cref{sec:related-work} positions the present work relative to the literature. \Cref{sec:repeated-composition} develops the discrete stability and error-propagation analysis. Sections~\ref{sec:latent-architecture}--\ref{sec:neural-param} construct the latent architecture and its neural parametrization. \Cref{sec:benchmark} documents numerical experiments. \Cref{sec:boundary-semiflow} derives the boundary-admissible finite-dimensional theory and the PDE-dependent theorem obligations. Sections~\ref{sec:transfer}--\ref{sec:expressivity} give the conditional transfer and architecture-compatible component bounds.

\section{Relation to existing work}\label{sec:related-work}

The present formulation is connected to several lines of work without coinciding with any single one of them. We organize the discussion around five themes.

\paragraph{Nonlinear semigroup theory.}
The mathematical foundation for autonomous evolution PDEs is the theory of nonlinear semigroups \cite{Pazy1983Semigroups,CazenaveHaraux1998Semilinear,Barbu2010NonlinearDE}. The standard setup associates a closed operator $\cF$ on a Banach space $X$ with a solution semigroup $(\cS_t)_{t\ge 0}$ satisfying $\cS_0=\mathrm{Id}$ and $\cS_{t+s}=\cS_t\circ \cS_s$. Structural laws such as maximum principles and Lyapunov dissipation are classical features of this framework. Our paper takes these structural properties as the target specification for the learned model, rather than treating them as optional training penalties.

\paragraph{Structure-preserving numerical methods.}
The idea that numerical discretizations should inherit structural properties of the continuous PDE has a long history in numerical analysis. Invariant-region preserving schemes guarantee that the discrete solution remains in a physically admissible set \cite{KarlsenRisebro2003Invariant,HundsdorferVerwer2003Numerical,ChernogorovaTemov2022InvariantRegion}. Strong stability preserving (SSP) time integrators \cite{GottliebShuTadmor2001StrongStability} preserve convex constraints and total variation bounds. Geometric numerical integration \cite{HairerLubichWanner2006Geometric,SanzSerna2016Geometric} targets symplectic, volume-preserving, or energy-conserving properties. Implicit--explicit (IMEX) methods \cite{PareschiRusso2005IMEX} separate stiff from non-stiff terms to maintain stability. Our approach differs in that the structural guarantees are built into the architecture of the learned evolution family, rather than into a fixed numerical scheme applied to a known vector field.

\paragraph{Deep learning for PDEs.}
The use of neural networks for PDE problems was catalyzed by the dynamical systems perspective of E \cite{E2017Proposal} and by physics-informed neural networks (PINNs) \cite{RaissiPerdikarisKarniadakis2019PINNs}. Early work on neural-network-based ODE and PDE solvers \cite{LagarisFotiadis1998} demonstrated the feasibility of the approach, and high-dimensional PDE solvers based on deep learning \cite{HanJentzenE2018} extended the scope to previously intractable problems. Deep Galerkin methods \cite{SirignanoSpiliopoulos2018DGM} and neural-network-based PDE solvers \cite{RuthottoHaber2020DeepPDE} have since expanded the scope of neural approaches. Our work differs from these in target: rather than solving a single PDE instance, we seek to learn the solution operator family itself, as in operator learning approaches.

\paragraph{Operator learning and neural operators.}
DeepONet \cite{LuJinKarniadakis2021DeepONet,LuLu2020DeepONet} and the Fourier neural operator (FNO) \cite{LiKovachkiAzizzadenesheliEtAl2021FNO} approximate input--output maps between function spaces. The general neural operator framework \cite{KovachkiLiLiuEtAl2023NeuralOperator} provides universality results for such approximations. These approaches typically learn a single-step map; our emphasis is on a whole evolution family intended for repeated composition in time, with structural guarantees that persist under iteration.

\paragraph{Structure-preserving neural ODEs.}
Neural ODEs \cite{ChenRubanovaBettencourtDuvenaud2018NODE,KidgerMorrilFoster2020NeuralDE,DupontDoucetTeh2019} parametrize continuous-time dynamics with neural networks. Hamiltonian neural networks (HNNs) \cite{GreydanusDzambaYosinski2019HNN} and SympNets \cite{JinZhangZhuTangKarniadakis2020SympNets} encode symplectic structure. OnsagerNet \cite{YuTianELi2021OnsagerNet} enforces dissipative structure via the Onsager principle. Port-Hamiltonian neural networks \cite{Desjardins2024PHNN} encode dissipative structure through interconnection and damping, and Lyapunov-based neural approaches construct stability certificates for learned dynamics. Symplectic recurrent neural networks \cite{ZhongDeySusmelj2020SympRNN} and stable neural flows \cite{MassaroliPoliPark2020StableFlows} address long-time stability of learned dynamics. Invariant-region-preserving neural networks \cite{LiuGomezSuli2022InvariantRegion} are the closest in spirit to the present paper. Recent work on the interplay between discretization and regularization in neural ODEs \cite{GeshkovskiRigollet2025NeuralODE} further highlights the importance of understanding how learned dynamics behave under repeated composition. Our contribution adds the semigroup-level formulation and the explicit error-propagation analysis under repeated composition, while the architectural construction (sigmoid logit map + gradient flow with PSD mobility) shares the general principle of encoding constraints via diffeomorphisms to unconstrained spaces.

The distinction of the present paper is that it uses the semigroup viewpoint to keep time composition, invariant-region preservation, and energy decay in the same target specification, while leaving continuous-to-discrete PDE transfer as a separate analytical issue.

\section{Target class of learned semigroups}\label{sec:target-class}
We next specify the class of learned families that is admissible for this objective. The basic mathematical object is a parameterized family of learned operators
\begin{equation}\label{eq:learned-family}
\PhiT(\cdot,\tau) : \cK \to \cK, \qquad \tau\ge 0,
\end{equation}
intended to approximate the solution operator \(\cS_{\tau}\). The five requirements below fall into two groups: \eqref{eq:s1}--\eqref{eq:s3} concern approximation of the evolution family itself, whereas \eqref{eq:s4}--\eqref{eq:s5} guarantee that the learned repeated composition stays inside the admissible set and continues to dissipate the chosen Lyapunov functional.

\begin{definition}[Structure-preserving deep semigroup learner]\label{def:main-learner}
A family \(\PhiT(\cdot,\tau):\cK\to \cK\) is called a structure-preserving deep semigroup learner if it satisfies the following properties.
\begin{enumerate}[label=(S\arabic*)]
    \item \textbf{Identity at zero time:}
    \begin{equation}\label{eq:s1}
    \PhiT(u,0)=u, \qquad u\in \cK.
    \end{equation}

    \item \textbf{Generator consistency:}
    for \(u\in D(\cF)\cap \cK\),
    \begin{equation}\label{eq:s2}
    \PhiT(u,\tau)=u+\tau\cF(u)+o(\tau), \qquad \tau\downarrow 0.
    \end{equation}

    \item \textbf{Semigroup defect control:}
    there exists a defect function \(\eta_{\mathrm{sg}}:[0,\infty)^2\to [0,\infty)\) such that
    \begin{equation}\label{eq:s3}
    \norm{\PhiT(\PhiT(u,\tau_1),\tau_2)-\PhiT(u,\tau_1+\tau_2)} \le \eta_{\mathrm{sg}}(\tau_1,\tau_2),
    \end{equation}
    for all \(u\in \cK\) and all \(\tau_1,\tau_2\ge 0\), with \(\eta_{\mathrm{sg}}(\tau_1,\tau_2)\to 0\) as \((\tau_1,\tau_2)\to (0,0)\).

    \item \textbf{Maximum-principle preservation:}
    \begin{equation}\label{eq:s4}
    u\in \cK \Longrightarrow \PhiT(u,\tau)\in \cK, \qquad \tau\ge 0.
    \end{equation}

    \item \textbf{Dissipation preservation:}
    \begin{equation}\label{eq:s5}
    \cE(\PhiT(u,\tau)) \le \cE(u), \qquad u\in \cK, \qquad \tau\ge 0.
    \end{equation}
\end{enumerate}
\end{definition}

Definition~\ref{def:main-learner} isolates the class of learned operator families that are admissible for the present problem. The first three clauses encode approximation of the PDE semigroup, whereas \eqref{eq:s4} and \eqref{eq:s5} encode preservation of the maximum principle and of the dissipation law. In particular, \eqref{eq:s3} measures whether the learned family behaves coherently across different time increments, while \eqref{eq:s4}--\eqref{eq:s5} are the structural conditions that prevent rollout from leaving the physically meaningful regime. This definition therefore serves as the interface between the continuous semigroup problem and the architectural constructions introduced later. It should be read as a PDE-level target specification; the latent architecture constructed later verifies only a finite-dimensional structural subset of these requirements, while generator matching and approximation to the target PDE semigroup remain separate issues.

\begin{remark}[Finite-dimensional analogue]\label{rem:finite-dim-def}
For the semidiscrete model considered later, the PDE-level Definition~\ref{def:main-learner} has a finite-dimensional counterpart. Let $X_N=\R^N$, $\cK_N=[m,M]^N$, and let $\cFN:\cK_N\to \R^N$ be a semidiscrete generator. A family $\PhiT(\cdot,\tau):\cK_N\to \cK_N$ is a \emph{finite-dimensional structure-preserving learner} if it satisfies \eqref{eq:s1} on $\cK_N$, the finite-dimensional generator relation
\begin{equation}\label{eq:s2-fd}
\PhiT(u,\tau)=u+\tau\cFN(u)+o(\tau), \qquad u\in \cK_N, \qquad \tau\downarrow 0,
\end{equation}
and \eqref{eq:s3}--\eqref{eq:s5} with $\cK$ replaced by $\cK_N$ and $\cE$ replaced by a discrete energy $\cE_N:\cK_N\to \R$. The latent architecture constructed in Sections~\ref{sec:latent-architecture}--\ref{sec:neural-param} verifies this finite-dimensional version on the interior set $(m,M)^N$, as detailed in Corollary~\ref{cor:verification}. The passage from the finite-dimensional learner back to the PDE-level specification is the subject of Problem~\ref{prob:transfer} in Section~\ref{sec:analytical-questions}.
\end{remark}

\section{Repeated composition and discrete stability}\label{sec:repeated-composition}
Fix a time step \(\Delta t>0\). Once the family \(\PhiT\) is learned, it generates a discrete solver by repeated composition:
\begin{equation}\label{eq:discrete-solver}
u^{n+1}=\PhiT(u^n,\Delta t), \qquad u^0=u_0,
\end{equation}
or equivalently
\begin{equation}\label{eq:iterated-map}
u^n=\bigl(\PhiT(\cdot,\Delta t)\bigr)^{\circ n}(u_0), \qquad n\in \N\cup\{0\}.
\end{equation}

This section clarifies what is gained by learning the semigroup rather than only a one-step map. Once the model is used for repeated time-stepping, two questions become central: does the discrete orbit remain in the admissible set and dissipate the Lyapunov functional, and how do one-step approximation error and semigroup inconsistency accumulate over many compositions?

\begin{theorem}[Discrete stability inherited from structure]\label{thm:stability}
Assume that \eqref{eq:s4} and \eqref{eq:s5} hold. Then, for every \(u_0\in \cK\), the discrete orbit \eqref{eq:iterated-map} satisfies
\begin{equation}\label{eq:discrete-stability}
u^n\in \cK, \qquad \cE(u^{n+1})\le \cE(u^n)\le \cE(u^0), \qquad n\in \N\cup\{0\}.
\end{equation}
In particular, the learned solver preserves the admissible set and satisfies a discrete Lyapunov law.
\end{theorem}

\begin{proof}
The proof is an induction on \(n\). Since \(u^0=u_0\in \cK\), the base step is immediate. If \(u^n\in \cK\), then \eqref{eq:s4} gives \(u^{n+1}=\PhiT(u^n,\Delta t)\in \cK\), and \eqref{eq:s5} yields
\begin{equation*}
\cE(u^{n+1})=\cE(\PhiT(u^n,\Delta t))\le \cE(u^n).
\end{equation*}
This proves \eqref{eq:discrete-stability}.
\end{proof}

Theorem~\ref{thm:stability} shows that the structural constraints are not merely formal conditions: they immediately turn the learned operator into a stable time-marching scheme. In particular, once \eqref{eq:s4} and \eqref{eq:s5} are built into the model, admissibility and Lyapunov decay propagate automatically to every iterate.

\begin{proposition}[Error decomposition through semigroup defect]\label{prop:error-decomposition}
Let \(T>0\) and \(u_0\in \cK\), and assume that \eqref{eq:s3} and \eqref{eq:s4} hold. Assume also that there exists \(L_T\ge 0\) such that
\begin{equation}\label{eq:lipschitz-step}
\norm{\PhiT(v,\tau)-\PhiT(w,\tau)} \le (1+L_T\tau)\norm{v-w}, \qquad 0\le \tau\le T,
\end{equation}
for all \(v,w\in \cK\). Define
\begin{equation}\label{eq:app-defect}
\varepsilon_{\mathrm{app}}(T;u_0):=\sup_{0\le t\le T}\norm{\PhiT(u_0,t)-\cS_tu_0},
\end{equation}
and
\begin{equation}\label{eq:sg-defect-dt}
\varepsilon_{\mathrm{sg}}(T,\Delta t):=\sup_{0\le t\le T-\Delta t}\eta_{\mathrm{sg}}(t,\Delta t).
\end{equation}
Then, for every integer \(n\) with \(0\le n\Delta t\le T\),
\begin{equation}\label{eq:error-bound}
\norm{\bigl(\PhiT(\cdot,\Delta t)\bigr)^{\circ n}(u_0)-\cS_{n\Delta t}u_0}
\le e^{L_TT}n\,\varepsilon_{\mathrm{sg}}(T,\Delta t)+\varepsilon_{\mathrm{app}}(T;u_0).
\end{equation}
\end{proposition}

\begin{proof}
Set
\begin{equation}\label{eq:delta-def}
\delta_n:=\norm{\bigl(\PhiT(\cdot,\Delta t)\bigr)^{\circ n}(u_0)-\PhiT(u_0,n\Delta t)}.
\end{equation}
Then \(\delta_0=0\). By \eqref{eq:s4}, the iterated state \((\PhiT(\cdot,\Delta t))^{\circ n}(u_0)\) remains in \(\cK\); moreover \(\PhiT(u_0,n\Delta t)\in \cK\). Hence the Lipschitz estimate applies to the two arguments below. For \(n\Delta t\le T-\Delta t\),
\begin{align*}
\delta_{n+1}
&\le \norm{\PhiT\bigl((\PhiT(\cdot,\Delta t))^{\circ n}(u_0),\Delta t\bigr)-\PhiT\bigl(\PhiT(u_0,n\Delta t),\Delta t\bigr)} \\
&\quad + \norm{\PhiT\bigl(\PhiT(u_0,n\Delta t),\Delta t\bigr)-\PhiT(u_0,(n+1)\Delta t)} \\
&\le (1+L_T\Delta t)\delta_n+\varepsilon_{\mathrm{sg}}(T,\Delta t).
\end{align*}
A discrete Gronwall argument gives the displayed bound below; we use the coarse estimate \(\sum_{j=0}^{n-1}(1+L_T\Delta t)^j\le n e^{L_TT}\), which is sufficient for the convergence statement.
\begin{equation}\label{eq:delta-bound}
\delta_n\le e^{L_TT}n\,\varepsilon_{\mathrm{sg}}(T,\Delta t).
\end{equation}
Combining \eqref{eq:delta-bound} with the triangle inequality yields \eqref{eq:error-bound}.
\end{proof}

Proposition~\ref{prop:error-decomposition} separates two conceptually different error sources. The term \(\varepsilon_{\mathrm{app}}\) measures the trajectory-level semigroup fit from the initial state \(u_0\), whereas \(\varepsilon_{\mathrm{sg}}\) measures the internal inconsistency of the learned composition law. The estimate shows that long-time repeated composition is controlled not only by short-time approximation quality, but also by how well the learned family behaves as a semigroup. This bound is conditional on the stability hypothesis \eqref{eq:lipschitz-step} and does not by itself provide a PDE-level approximation theorem.

\begin{remark}[Domain of the Lipschitz condition]\label{rem:lipschitz-domain}
Proposition~\ref{prop:error-decomposition} assumes \eqref{eq:lipschitz-step} holds on the full admissible set $\cK$. In practice, the Lipschitz verification (Proposition~\ref{prop:lipschitz} below) establishes the condition only on compact subsets $\cU\subset(m,M)^N$, because $g^{-1}$ is Lipschitz only on compact subsets of the interior. The error bound \eqref{eq:error-bound} therefore applies to initial data $u_0$ whose trajectory under $\PhiT$ remains in a compact subset of the interior for $0\le t\le T$. This restriction is natural: states near the boundary $\{u_i=m\text{ or }u_i=M\}$ have $g^{-1}(u_i)\to\pm\infty$, and the latent dynamics may degenerate there.
\end{remark}

\begin{corollary}[A route to convergence]\label{cor:global-convergence}
Let \(T>0\) and \(u_0\in \cK\). Assume that, for a sequence \(\Delta t_m\downarrow 0\) and corresponding learned models satisfying \eqref{eq:s3}--\eqref{eq:s4}, one has
\begin{equation}\label{eq:asymptotic-defects}
\varepsilon_{\mathrm{app}}^{(m)}(T;u_0)\to 0, \qquad \frac{\varepsilon_{\mathrm{sg}}^{(m)}(T,\Delta t_m)}{\Delta t_m}\to 0,
\end{equation}
and that \eqref{eq:lipschitz-step} holds with a constant \(L_T\) independent of \(m\). Then
\begin{equation}\label{eq:asymptotic-conv}
\sup_{0\le n\Delta t_m\le T}\norm{\bigl(\PhiT^{(m)}(\cdot,\Delta t_m)\bigr)^{\circ n}(u_0)-\cS_{n\Delta t_m}u_0}\to 0
\end{equation}
as \(m\to \infty\).
\end{corollary}

\begin{proof}
If \(0\le n\Delta t_m\le T\), then \(n\le T/\Delta t_m\). Proposition~\ref{prop:error-decomposition} gives
\begin{equation*}
\norm{\bigl(\PhiT^{(m)}(\cdot,\Delta t_m)\bigr)^{\circ n}(u_0)-\cS_{n\Delta t_m}u_0}
\le e^{L_TT}T\frac{\varepsilon_{\mathrm{sg}}^{(m)}(T,\Delta t_m)}{\Delta t_m}+\varepsilon_{\mathrm{app}}^{(m)}(T;u_0).
\end{equation*}
Taking the supremum over \(n\) and using \eqref{eq:asymptotic-defects} proves \eqref{eq:asymptotic-conv}.
\end{proof}

\section{A latent structure-preserving architecture}\label{sec:latent-architecture}
Corollary~\ref{cor:global-convergence} therefore suggests a concrete design principle: a learned evolution family should control both approximation error and semigroup defect, while preserving the structural laws that guarantee stable repeated composition. The next question is how to build these requirements into an architecture for which the structural clauses of Definition~\ref{def:main-learner} hold at the model-class level rather than only through soft penalties during training.

We now work with a semidiscrete finite-dimensional state model (Assumption~\ref{ass:finite-dim} below). The exact statements proved in the remainder of the paper are therefore architecture-level finite-dimensional guarantees, not direct PDE-level transfer results. The rigorous passage from this exact finite-dimensional verification back to the target infinite-dimensional semigroup is deferred to Section~\ref{sec:analytical-questions}.

\begin{assumption}[Finite-dimensional state model]\label{ass:finite-dim}
Fix \(N\ge 1\), and identify the spatially discretized state space with
\begin{equation}\label{eq:XN}
X_N=\R^N.
\end{equation}
Assume that the maximum principle is represented by the box constraint
\begin{equation}\label{eq:KN}
\cK_N=[m,M]^N=\{u=(u_1,\dots,u_N)^\top\in \R^N : m\le u_i\le M, \ i=1,\dots,N\}.
\end{equation}
\end{assumption}

We regard \(X_N\) and \(\cK_N\) as the finite-dimensional counterparts of \(X\) and \(\cK\), respectively. For the inverse-logit latent construction used below, the actual input domain is the interior set \((m,M)^N\), because the latent inverse \(g^{-1}\) is only defined there. The resulting learned family is therefore considered as a map
\begin{equation}\label{eq:finite-learner}
\PhiT : (m,M)^N\times [0,\infty)\to \cK_N.
\end{equation}

The construction below follows a simple architectural principle: first lift an interior constrained state into an unconstrained latent variable, then evolve that latent variable by an autonomous dissipative flow, and finally map the latent state back into the admissible set. Boundary states with some component equal to \(m\) or \(M\) are not covered by the present parametrization.

\begin{definition}[Bounded latent parametrization]\label{def:latent-map}
Let \(\sigma:\R\to (0,1)\) be the sigmoid function
\begin{equation}\label{eq:sigmoid}
\sigma(\xi)=\frac{1}{1+e^{-\xi}}.
\end{equation}
Define \(g:\R^N\to (m,M)^N\) componentwise by
\begin{equation}\label{eq:g-map}
g(z)_i=m+(M-m)\sigma(z_i), \qquad i=1,\dots,N.
\end{equation}
Then \(g\) maps the whole latent space into the admissible set. Its inverse on \((m,M)^N\) is
\begin{equation}\label{eq:g-inv}
g^{-1}(u)_i=\log\!\left(\frac{u_i-m}{M-u_i}\right), \qquad i=1,\dots,N,
\end{equation}
and its Jacobian is
\begin{equation}\label{eq:Dg}
\mathrm{D}g(z)=(M-m)\,\diag\bigl(\sigma(z_i)(1-\sigma(z_i))\bigr)_{i=1}^N.
\end{equation}
\end{definition}

\begin{definition}[Latent dissipative semigroup]\label{def:latent-semigroup}
Assume that
\begin{equation}\label{eq:latent-objects}
\Psi_{\theta}\in C^1(\R^N;\R), \qquad K_{\theta}:\R^N\to \R^{N\times N},
\end{equation}
and that \(K_{\theta}(z)\) is symmetric positive semidefinite for every \(z\in \R^N\). Consider the latent ODE
\begin{equation}\label{eq:latent-ode}
\partial_s Z_{\theta}(s;z_0)=-K_{\theta}(Z_{\theta}(s;z_0))\nabla \Psi_{\theta}(Z_{\theta}(s;z_0)), \qquad Z_{\theta}(0;z_0)=z_0.
\end{equation}
Assume furthermore that this vector field generates a local flow \(\varphi_{\theta}^{\tau}\). The learned operator family is then defined for interior states and for those times \(\tau\ge 0\) for which the latent trajectory through \(g^{-1}(u)\) exists, by
\begin{equation}\label{eq:latent-learner}
\PhiT(u,\tau):=g\bigl(\varphi_{\theta}^{\tau}(g^{-1}(u))\bigr).
\end{equation}
The induced discrete energy on $(m,M)^N$ is
\begin{equation}\label{eq:induced-energy}
\cE_{\theta}(u):=\Psi_{\theta}(g^{-1}(u)).
\end{equation}
\end{definition}

\begin{remark}[Learned, discrete, and physical energies]\label{rem:energy-distinction}
The induced functional $\cE_\theta$ is a learned Lyapunov function. It is not automatically equal to a discrete physical energy $\cE_N$, nor does its monotonicity imply monotonicity of the PDE functional $\cE$. Establishing $\cE_\theta\approx\cE_N$, or separately evaluating both functionals, is an additional requirement. Accordingly, the experiments and code use distinct labels for learned-energy and physical-energy monotonicity.
\end{remark}

\begin{remark}[Regularity hierarchy]\label{rem:regularity}
Definition~\ref{def:latent-semigroup} assumes $\Psi_{\theta}\in C^1$, which suffices for the structural properties (Proposition~\ref{prop:structure}). The global-flow existence result (Proposition~\ref{prop:global-flow}) and the Lipschitz verification (Proposition~\ref{prop:lipschitz}) require the stronger regularity $\Psi_{\theta}\in C^2$ and $K_{\theta}\in C^1$, which is ensured by the neural parametrization of Section~\ref{sec:neural-param} when the activation $\rho$ is $C^2$ (e.g., softplus or tanh).
\end{remark}

\begin{proposition}[Sufficient conditions for global flow]\label{prop:global-flow}
Assume that $\Psi_{\theta}\in C^2(\R^N;\R)$ satisfies the coercivity condition
\begin{equation}\label{eq:coercivity}
\Psi_{\theta}(z)\ge c_1\norm{z}^2-c_2, \qquad z\in \R^N,
\end{equation}
for some constants $c_1>0$ and $c_2\ge 0$. Assume furthermore that $K_{\theta}\in C^1(\R^N;\R^{N\times N})$ is symmetric positive semidefinite for every $z\in \R^N$. Then the latent ODE \eqref{eq:latent-ode} generates a global flow on $\R^N$.
\end{proposition}

\begin{proof}
Let $z_0\in \R^N$ and let $Z_{\theta}(s;z_0)$ denote the maximal solution. Since $\Psi_{\theta}\in C^2$ and $K_{\theta}\in C^1$, the vector field $f_{\theta}(z)=-K_{\theta}(z)\nabla \Psi_{\theta}(z)$ is of class $C^1$, hence locally Lipschitz on $\R^N$. The energy dissipation identity gives
\begin{equation*}
\frac{d}{ds}\Psi_{\theta}(Z_{\theta}(s))=-\nabla \Psi_{\theta}(Z_{\theta}(s))^{\top}K_{\theta}(Z_{\theta}(s))\nabla \Psi_{\theta}(Z_{\theta}(s))\le 0,
\end{equation*}
so $\Psi_{\theta}(Z_{\theta}(s))\le \Psi_{\theta}(z_0)$ for all $s$ in the maximal interval of existence. By coercivity \eqref{eq:coercivity}, $\norm{Z_{\theta}(s)}^2\le (\Psi_{\theta}(z_0)+c_2)/c_1$, so the solution remains bounded on any finite interval. By the extension theorem for ODEs, the maximal solution exists for all $s\ge 0$, i.e., the flow is global.
\end{proof}

\begin{remark}[Verifiability of the global-flow conditions]\label{rem:global-flow-verify}
For the neural parametrization of Section~\ref{sec:neural-param}, condition \eqref{eq:coercivity} holds if $\beta_V>0$ in \eqref{eq:V-mlp}, since the quadratic term $\frac{\beta_V}{2}\norm{z}^2$ dominates for large $\norm{z}$. The regularity $K_{\theta}\in C^1$ and positive semidefiniteness hold automatically for the softplus-based mobility parametrization of \eqref{eq:k-param}. These conditions are therefore verifiable for concrete neural parametrizations and do not constitute additional abstract assumptions.
\end{remark}

\begin{remark}[Archived noncoercive configuration and revised option]\label{rem:betaV-zero}
The archived checkpoints reported in Section~\ref{sec:benchmark} were initialized with $\beta_V=0$. They are therefore not covered by Proposition~\ref{prop:global-flow}; absence of numerical blowup on finitely many trajectories is only a diagnostic and is not evidence of global existence. The revised implementation adds a non-trainable coefficient floor $\beta_{V,\min}\ge0$ and uses the effective coefficient $\beta_V^{\mathrm{eff}}=\beta_{V,\min}+|\widetilde\beta_V|$. Selecting $\beta_{V,\min}>0$ prevents optimization from removing the quadratic term and makes the coercivity configuration explicit. Results from such a configuration must be reported separately rather than attributed retroactively to archived checkpoints.
\end{remark}

\begin{proposition}[Explicit global-flow criterion for the neural parametrization]\label{prop:explicit-global-flow}
Consider the neural parametrization of Definitions~\ref{def:rd-energy}--\ref{def:neural-param}. Suppose the following hold:
\begin{enumerate}[label=(\roman*)]
    \item The quadratic coefficient satisfies $\beta_V>0$ in \eqref{eq:V-mlp}.
    \item The MLP part $V_{\theta}^{\mathrm{mlp}}(\xi):=(c^V)^\top h_{L_V-1}^V(\xi)$ satisfies a growth bound $|V_{\theta}^{\mathrm{mlp}}(\xi)|\le C_1(1+|\xi|)$ for some $C_1>0$.
    \item The mobility coefficients are uniformly bounded: $k_{\theta,i}(z)\le K_{\max}$ for all $z\in \R^N$ and $i=1,\dots,N$, where $K_{\max}<\infty$ is a constant depending only on the network parameters.
\end{enumerate}
Then the latent ODE \eqref{eq:latent-ode} generates a global flow on $\R^N$.
\end{proposition}

\begin{proof}
We verify the hypotheses of Proposition~\ref{prop:global-flow}. Condition~(i) ensures $\Psi_{\theta}(z)\ge \frac{\beta_V}{2}\norm{z}^2 - C_2$ for some $C_2\ge 0$, which is the coercivity condition \eqref{eq:coercivity} with $c_1=\beta_V/2$. Condition~(ii) ensures the MLP component does not overwhelm the quadratic term at large $\norm{z}$, preserving the coercivity of $\Psi_{\theta}$. The mobility $K_{\theta}$ is symmetric positive semidefinite by construction (softplus output). The $C^1$ regularity of $K_{\theta}$ follows from the smooth activation $\rho=\softplus$. The result then follows from Proposition~\ref{prop:global-flow}.
\end{proof}

\begin{proposition}[Structural properties of the latent-flow construction]\label{prop:structure}
Assume that \eqref{eq:latent-ode} generates a global flow on \(\R^N\). Then, for every \(u\in (m,M)^N\) and every \(\tau,\tau_1,\tau_2\ge 0\), the map \eqref{eq:latent-learner} is well defined, takes values in \((m,M)^N\), and satisfies
\begin{equation}\label{eq:structure-p1}
\PhiT(u,0)=u,
\end{equation}
\begin{equation}\label{eq:structure-p3}
\PhiT(\PhiT(u,\tau_1),\tau_2)=\PhiT(u,\tau_1+\tau_2),
\end{equation}
\begin{equation}\label{eq:structure-p4}
\PhiT(u,\tau)\in \cK_N,
\end{equation}
and
\begin{equation}\label{eq:structure-p5}
\cE_{\theta}(\PhiT(u,\tau))\le \cE_{\theta}(u).
\end{equation}
\end{proposition}

\begin{proof}
Set \(z_0:=g^{-1}(u)\). Equation \eqref{eq:structure-p1} follows from the identity property of the flow. Equation \eqref{eq:structure-p3} follows from the semigroup property of the autonomous flow \(\varphi_{\theta}^{\tau}\). Since \(g(\R^N)\subset (m,M)^N\subset \cK_N\), the map \eqref{eq:latent-learner} is well defined for all \(\tau\ge 0\), takes values in \((m,M)^N\), and in particular satisfies equation \eqref{eq:structure-p4}. Finally,
\begin{align*}
\frac{d}{ds}\Psi_{\theta}(Z_{\theta}(s;z_0))
&=\nabla \Psi_{\theta}(Z_{\theta}(s;z_0))^\top \partial_s Z_{\theta}(s;z_0) \\
&=-\nabla \Psi_{\theta}(Z_{\theta}(s;z_0))^\top K_{\theta}(Z_{\theta}(s;z_0))\nabla \Psi_{\theta}(Z_{\theta}(s;z_0)) \\
&\le 0.
\end{align*}
Hence \eqref{eq:structure-p5} follows from \eqref{eq:induced-energy}.
\end{proof}

\begin{proposition}[Generator expansion and matching condition]\label{prop:generator}
Assume that \(\Psi_{\theta}\in C^2(\R^N)\) and \(K_{\theta}\in C^1(\R^N;\R^{N\times N})\). Then the latent vector field in \eqref{eq:latent-ode} is of class \(C^1\), so it generates a local \(C^1\) flow. For each fixed \(u\in (m,M)^N\),
\begin{equation}\label{eq:generator-expansion}
\PhiT(u,\tau)=u-\tau\,\mathrm{D}g(g^{-1}(u))K_{\theta}(g^{-1}(u))\nabla \Psi_{\theta}(g^{-1}(u))+o(\tau), \qquad \tau\downarrow 0.
\end{equation}
Consequently, if there exists a semidiscrete generator \(\cFN:(m,M)^N\to \R^N\) such that the matching condition
\begin{equation}\label{eq:matching}
\mathrm{D}g(g^{-1}(u))K_{\theta}(g^{-1}(u))\nabla \Psi_{\theta}(g^{-1}(u))=-\cFN(u), \qquad u\in (m,M)^N,
\end{equation}
is satisfied, then
\begin{equation*}
\PhiT(u,\tau)=u+\tau\cFN(u)+o(\tau), \qquad \tau\downarrow 0,
\end{equation*}
which is the finite-dimensional counterpart of \eqref{eq:s2}.
\end{proposition}

\begin{proof}
Fix \(u\in (m,M)^N\), and set \(z_0:=g^{-1}(u)\). Since the latent vector field
\begin{equation*}
f_{\theta}(z):=-K_{\theta}(z)\nabla \Psi_{\theta}(z)
\end{equation*}
is \(C^1\), the local flow satisfies the standard first-order expansion
\begin{equation*}
\varphi_{\theta}^{\tau}(z_0)=z_0+\tau f_{\theta}(z_0)+o(\tau), \qquad \tau\downarrow 0.
\end{equation*}
Because \(g\in C^1(\R^N;(m,M)^N)\), Taylor expansion at \(z_0\) gives
\begin{align*}
\PhiT(u,\tau)
&=g\bigl(\varphi_{\theta}^{\tau}(z_0)\bigr) \\
&=g(z_0)+\mathrm{D}g(z_0)\bigl(\varphi_{\theta}^{\tau}(z_0)-z_0\bigr)+o\!\left(\norm{\varphi_{\theta}^{\tau}(z_0)-z_0}\right) \\
&=u-\tau\,\mathrm{D}g(z_0)K_{\theta}(z_0)\nabla\Psi_{\theta}(z_0)+o(\tau).
\end{align*}
This is \eqref{eq:generator-expansion}. If \eqref{eq:matching} holds for some semidiscrete generator \(\cFN\), then
\begin{equation*}
\PhiT(u,\tau)=u+\tau\cFN(u)+o(\tau), \qquad \tau\downarrow 0,
\end{equation*}
which is the finite-dimensional counterpart of \eqref{eq:s2}.
\end{proof}

\begin{corollary}[Verification of the finite-dimensional structural clauses]\label{cor:verification}
Assume that the hypotheses of Proposition~\ref{prop:structure} hold. Then the map \eqref{eq:latent-learner} satisfies \eqref{eq:s1} and \eqref{eq:s4} exactly on the interior state set \((m,M)^N\), it satisfies \eqref{eq:s3} there with defect function
\begin{equation}\label{eq:exact-defect}
\eta_{\mathrm{sg}}(\tau_1,\tau_2)\equiv 0.
\end{equation}
and it satisfies the finite-dimensional dissipation law
\begin{equation}\label{eq:exact-discrete-energy}
\cE_{\theta}(\PhiT(u,\tau))\le \cE_{\theta}(u), \qquad u\in (m,M)^N, \qquad \tau\ge 0.
\end{equation}
If, in addition, Proposition~\ref{prop:generator} holds and
\begin{equation}\label{eq:matching-exact}
\mathrm{D}g(g^{-1}(u))K_{\theta}(g^{-1}(u))\nabla \Psi_{\theta}(g^{-1}(u))=-\cFN(u), \qquad u\in (m,M)^N,
\end{equation}
then the finite-dimensional generator relation
\begin{equation}\label{eq:s2-finite}
\PhiT(u,\tau)=u+\tau\cFN(u)+o(\tau), \qquad u\in (m,M)^N,
\end{equation}
holds as \(\tau\downarrow 0\).
\end{corollary}

Consequently, the latent-flow construction provides an exact finite-dimensional verification of the identity, admissibility, exact semigroup, and induced-energy dissipation properties on the interior state set, while generator matching remains conditional on \eqref{eq:matching-exact} and concerns a semidiscrete generator \(\cFN\), not yet the original PDE generator \(\cF\).

\begin{proof}
Equation \eqref{eq:s1} follows from \eqref{eq:structure-p1}; equation \eqref{eq:s4} follows from \eqref{eq:structure-p4}; \eqref{eq:s3} follows from \eqref{eq:structure-p3}, which implies \eqref{eq:exact-defect}; and \eqref{eq:exact-discrete-energy} follows from \eqref{eq:structure-p5}. The final assertion is immediate from Proposition~\ref{prop:generator} and \eqref{eq:matching-exact}.
\end{proof}

\begin{remark}[Continuous-time vs.\ discrete-time semigroup]\label{rem:cont-vs-disc}
The exact semigroup property \eqref{eq:exact-defect} ($\eta_{\mathrm{sg}}\equiv 0$) holds only for the exact continuous-time latent flow. A fourth-order Runge--Kutta implementation generally has a nonzero numerical composition defect. The archived value near $10^{-14}$ for $M_{\mathrm{sub}}=30$ is a floating-point observation for one protocol, not an exact guarantee and not by itself a verified integrator-error estimate. The revised evaluator reports absolute and relative numerical defects and supports substep refinement while restoring the model's original substep count.
\end{remark}

\section{Neural parametrization of the latent semigroup}\label{sec:neural-param}
The previous section identifies the abstract ingredients needed for a structure-preserving semigroup learner. We now turn those ingredients into explicit neural parametrizations of the latent energy, interaction matrix, and mobility.
\begin{definition}[Reaction--diffusion latent energy ansatz]\label{def:rd-energy}
Let
\begin{equation}\label{eq:V-type}
V_{\theta}\in C^1(\R;\R)
\end{equation}
be a scalar potential, and let
\begin{equation}\label{eq:A-sym}
a_{ij}^{(\theta)}=a_{ji}^{(\theta)}\ge 0, \qquad i,j=1,\dots,N.
\end{equation}
Define
\begin{equation}\label{eq:rd-energy}
\Psi_{\theta}(z)=\sum_{i=1}^N V_{\theta}(z_i)+\frac12\sum_{i,j=1}^N a_{ij}^{(\theta)}(z_i-z_j)^2.
\end{equation}
If, in addition,
\begin{equation}\label{eq:diag-mobility}
K_{\theta}(z)=\diag\bigl(k_{\theta,1}(z),\dots,k_{\theta,N}(z)\bigr), \qquad k_{\theta,i}(z)\ge 0,
\end{equation}
then the latent dynamics takes the componentwise form
\begin{equation}\label{eq:component-flow}
\partial_s z_i=-k_{\theta,i}(z)\left(V_{\theta}'(z_i)+2\sum_{j=1}^N a_{ij}^{(\theta)}(z_i-z_j)\right), \qquad i=1,\dots,N.
\end{equation}
\end{definition}

\begin{definition}[Neural parametrization of latent energy and mobility]\label{def:neural-param}
Fix an activation function \(\rho:\R\to \R\) of class \(C^1\), and define
\begin{equation}\label{eq:softplus}
\softplus(\xi)=\log(1+e^{\xi}).
\end{equation}
For the scalar potential, fix integers \(d_1^V,\dots,d_{L_V-1}^V\in \N\), together with parameters
\begin{equation}\label{eq:V-params}
W_0^V\in \R^{d_1^V\times 1}, \qquad b_0^V\in \R^{d_1^V},
\end{equation}
\begin{equation*}
W_{\ell}^V\in \R^{d_{\ell+1}^V\times d_{\ell}^V}, \qquad b_{\ell}^V\in \R^{d_{\ell+1}^V}, \quad \ell=1,\dots,L_V-2,
\qquad c^V\in \R^{d_{L_V-1}^V}, \qquad \beta_V\ge 0.
\end{equation*}
Define
\begin{equation}\label{eq:V-hidden}
h_1^V(\xi):=\rho\bigl(W_0^V\xi+b_0^V\bigr), \qquad
h_{\ell+1}^V(\xi):=\rho\bigl(W_{\ell}^V h_{\ell}^V(\xi)+b_{\ell}^V\bigr), \quad \ell=1,\dots,L_V-2,
\end{equation}
and set
\begin{equation}\label{eq:V-mlp}
V_{\theta}(\xi):=(c^V)^\top h_{L_V-1}^V(\xi)+\frac{\beta_V}{2}\xi^2.
\end{equation}

For the interaction coefficients, fix an integer \(r_A\ge 1\), choose symmetric binary weights
\(m_{ij}=m_{ji}\in \{0,1\}\) and embeddings \(e_i\in \R^{r_A}\) for \(i=1,\dots,N\), and define
\begin{equation}\label{eq:A-param}
a_{ij}^{(\theta)}:=m_{ij}\,\softplus(e_i^\top e_j), \qquad i,j=1,\dots,N.
\end{equation}

For the diagonal mobility, fix a stencil radius \(r_K\in \N\) and integers \(d_1^K,\dots,d_{L_K-1}^K\in \N\). Fix parameters
\begin{equation}\label{eq:K-params}
W_0^K\in \R^{d_1^K\times (2r_K+1)}, \qquad b_0^K\in \R^{d_1^K},
\end{equation}
\begin{equation*}
W_{\ell}^K\in \R^{d_{\ell+1}^K\times d_{\ell}^K}, \qquad b_{\ell}^K\in \R^{d_{\ell+1}^K}, \quad \ell=1,\dots,L_K-2,
\qquad c^K\in \R^{d_{L_K-1}^K}, \qquad \beta_K\in \R.
\end{equation*}
For each \(z\in\R^N\) and \(i=1,\dots,N\), extend boundary values by zero-padding (setting \(z_j=0\) for indices \(j\notin\{1,\dots,N\}\)) and define
\begin{equation}\label{eq:K-hidden}
h_{i,1}^K(z):=\rho\bigl(W_0^K(z_{i-r_K},\dots,z_i,\dots,z_{i+r_K})^\top+b_0^K\bigr),
\end{equation}
\begin{equation*}
h_{i,\ell+1}^K(z):=\rho\bigl(W_{\ell}^K h_{i,\ell}^K(z)+b_{\ell}^K\bigr), \quad \ell=1,\dots,L_K-2,
\end{equation*}
and set
\begin{equation}\label{eq:k-param}
k_{\theta,i}(z):=\softplus\bigl((c^K)^\top h_{i,L_K-1}^K(z)+\beta_K\bigr), \qquad i=1,\dots,N.
\end{equation}
Hence the diagonal mobility matrix is
\begin{equation}\label{eq:K-matrix}
K_{\theta}(z):=\diag\bigl(k_{\theta,1}(z),\dots,k_{\theta,N}(z)\bigr).
\end{equation}
\end{definition}

The zero-padding convention is used only to make the finite-dimensional stencil map well defined. It can be replaced by another fixed extension rule, such as periodic, reflective, or boundary-condition-aware padding, without changing the abstract latent-flow argument above.

\begin{proposition}[Parameterized latent dynamics]\label{prop:param-dynamics}
Under Definitions~\ref{def:rd-energy} and \ref{def:neural-param}, \(\Psi_\theta\in C^1(\R^N)\) and \(K_\theta(z)\) is symmetric positive semidefinite for every \(z\in\R^N\). These identities verify the regularity and positivity conditions in Definition~\ref{def:latent-semigroup}. Moreover, the coefficients in \eqref{eq:rd-energy} and \eqref{eq:component-flow} are realized by \eqref{eq:A-param} and \eqref{eq:k-param}. Consequently,
\begin{equation}\label{eq:psi-param}
\Psi_{\theta}(z)=\sum_{i=1}^N V_{\theta}(z_i)+\frac12\sum_{i,j=1}^N a_{ij}^{(\theta)}(z_i-z_j)^2,
\end{equation}
its gradient satisfies
\begin{equation}\label{eq:psi-grad}
\bigl(\nabla\Psi_{\theta}(z)\bigr)_i=(V_{\theta})'(z_i)+2\sum_{j=1}^N a_{ij}^{(\theta)}(z_i-z_j), \qquad i=1,\dots,N,
\end{equation}
and the latent ODE becomes
\begin{equation}\label{eq:param-flow}
\partial_s z_i=-k_{\theta,i}(z)\left((V_{\theta})'(z_i)+2\sum_{j=1}^N a_{ij}^{(\theta)}(z_i-z_j)\right),
\end{equation}
for \(i=1,\dots,N\).
\end{proposition}

\begin{proof}
By Definition~\ref{def:neural-param}, the scalar network \(V_{\theta}\) is a finite composition of affine maps with the \(C^1\) activation \(\rho\), followed by the smooth quadratic term \(\frac{\beta_V}{2}\xi^2\); hence \(V_{\theta}\in C^1(\R)\). Likewise, each coefficient \(a_{ij}^{(\theta)}=m_{ij}\,\softplus(e_i^\top e_j)\) is symmetric and nonnegative, and each mobility coefficient \(k_{\theta,i}(z)\) is nonnegative because \(\softplus\ge 0\). Therefore \(K_{\theta}(z)\) is diagonal with nonnegative diagonal entries, hence symmetric positive semidefinite for every \(z\in \R^N\).

Since \(\Psi_{\theta}\) in \eqref{eq:rd-energy} is a finite sum of \(C^1\) terms, it belongs to \(C^1(\R^N)\). Differentiating \eqref{eq:rd-energy} with respect to \(z_i\) and using the symmetry \(a_{ij}^{(\theta)}=a_{ji}^{(\theta)}\) yields
\begin{equation*}
\partial_{z_i}\!\left(\frac12\sum_{p,q=1}^N a_{pq}^{(\theta)}(z_p-z_q)^2\right)=2\sum_{j=1}^N a_{ij}^{(\theta)}(z_i-z_j),
\end{equation*}
which gives \eqref{eq:psi-grad}. Substituting \eqref{eq:psi-grad} and \eqref{eq:k-param} into the latent ODE \eqref{eq:latent-ode} gives \eqref{eq:param-flow}. The opening assertions verify the regularity and positivity conditions appearing in Definition~\ref{def:latent-semigroup}.
\end{proof}

\begin{proposition}[One-step Lipschitz condition for the latent architecture]\label{prop:lipschitz}
Assume that the hypotheses of Proposition~\ref{prop:global-flow} hold and that the latent ODE generates a global flow. Fix $T>0$ and a compact subset $\cU\subset (m,M)^N$. If $\Psi_{\theta}\in C^2(\R^N)$ and $K_{\theta}\in C^1(\R^N;\R^{N\times N})$, then the latent vector field $f_{\theta}(z)=-K_{\theta}(z)\nabla \Psi_{\theta}(z)$ is locally Lipschitz on $\R^N$, and the flow $\varphi_{\theta}^{\tau}$ is Lipschitz on the preimage $g^{-1}(\cU)$ with constant depending on $T$ and $\cU$. Consequently, there exists $L_T\ge 0$ such that
\begin{equation}\label{eq:lipschitz-step-verified}
\norm{\PhiT(v,\tau)-\PhiT(w,\tau)} \le (1+L_T\tau)\norm{v-w}, \qquad v,w\in \cU, \qquad 0\le \tau\le T.
\end{equation}
In particular, the hypothesis \eqref{eq:lipschitz-step} of Proposition~\ref{prop:error-decomposition} holds on compact subsets of the interior state set.
\end{proposition}

\begin{proof}
Since $\Psi_{\theta}\in C^2$ and $K_{\theta}\in C^1$, the vector field $f_{\theta}(z)=-K_{\theta}(z)\nabla \Psi_{\theta}(z)$ is $C^1$, hence locally Lipschitz. On the compact set $g^{-1}(\cU)$, let $L_f$ denote the Lipschitz constant of $f_{\theta}$ on a compact neighborhood. By the standard Lipschitz dependence of ODE flows on initial data (obtained via Gronwall's inequality applied to the variational equation), the flow satisfies
\begin{equation*}
\norm{\varphi_{\theta}^{\tau}(z_1)-\varphi_{\theta}^{\tau}(z_2)}\le e^{L_f\tau}\norm{z_1-z_2}, \qquad z_1,z_2\in g^{-1}(\cU), \qquad 0\le \tau\le T.
\end{equation*}
The sigmoid map $g$ is Lipschitz on $\R^N$ with constant $(M-m)/4$, and $g^{-1}$ is Lipschitz on compact subsets of $(m,M)^N$ with constant $4/((M-m)\delta)$ where $\delta=\min_{u\in \cU}\min_i\{u_i-m,M-u_i\}>0$. Composing the three Lipschitz bounds gives the stated estimate with $L_T$ depending on $L_f$, $T$, and the geometry of $\cU$.
\end{proof}

\begin{remark}[Tightness at $\tau=0$]\label{rem:tightness-tau0}
At $\tau=0$, the composition bound from the proof above gives $\norm{\PhiT(v,0)-\PhiT(w,0)}\le \norm{v-w}/\delta$, which overestimates the true Lipschitz constant of $\PhiT(\cdot,0)=\mathrm{Id}$ (which is exactly $1$). The tighter form $\norm{\PhiT(v,\tau)-\PhiT(w,\tau)}\le (1+L_T\tau)\norm{v-w}$ stated in \eqref{eq:lipschitz-step-verified} follows from the mean value theorem applied to the Jacobian $D_v\PhiT(v,\tau)$, which equals $I$ at $\tau=0$ and varies continuously with $\tau$. This form preserves the linear-in-$\tau$ growth used in the Gronwall argument of Proposition~\ref{prop:error-decomposition}.
\end{remark}

\section{Benchmark experiments}\label{sec:benchmark}

We illustrate how the abstract ingredients of Sections~\ref{sec:latent-architecture}--\ref{sec:neural-param} apply to Fisher--KPP and Allen--Cahn reaction--diffusion equations. Viscous Burgers includes a transport term outside the pure latent gradient-flow model class and is included only as an empirical stress test. The numerical tables in this snapshot are archival: after correction of the reference-time alignment, every comparative value must be rerun before it is treated as revised evidence.

\subsection{Semigroup-based training objectives}\label{subsec:training}
Suppose now that training data contain snapshots of the target evolution, namely triples $(u,\tau,\cS_{\tau}u)$ with states sampled from the interior set on which the latent construction is defined. The training objective is the one-step prediction loss:
\begin{equation}\label{eq:loss-step}
\mathcal L_{\mathrm{step}}=\E\,\norm{\PhiT(u,\tau)-\cS_{\tau}u}^2,
\end{equation}

\paragraph{Training strategy in the experiments.}
The archived Fisher--KPP runner used $\mathcal L_{\mathrm{step}}+0.1\mathcal L_{\mathrm{rollout}}+0.01\mathcal L_{\mathrm{learned\ energy}}$ for the latent model, so its accuracy cannot be described as obtained from the one-step loss alone. These auxiliary terms are not needed to define the exact continuous architecture, but they can affect optimization and empirical accuracy. The revised runner therefore records all loss weights and provides an architecture-only mode with both auxiliary weights set to zero. Learned-energy and physical-energy diagnostics are reported separately.

\subsection{Fisher--KPP equation}\label{subsec:fisher-kpp}
Take the Fisher--KPP equation
\begin{equation}\label{eq:fisher-kpp-bench}
\partial_t u = \nu \Delta u + r u(1-u), \qquad 0\le u\le 1,
\end{equation}
on a spatial domain $\Omega\subset\R^d$ with homogeneous Neumann or periodic boundary conditions. The admissible set is $\cK=\{u:0\le u\le 1\}$ and the dynamics is the $L^2$-gradient flow of the energy
\begin{equation}\label{eq:fisher-energy}
\cE(u)=\int_\Omega\Bigl(\frac{\nu}{2}|\nabla u|^2 - r\Bigl(\frac{u^2}{2}-\frac{u^3}{3}\Bigr)\Bigr)\,dx,
\end{equation}
so that $\cE(\cS_t u)\le\cE(u)$ for all $t\ge 0$. The energy $\cE$ is bounded from below on $\cK$ since $r(u^2/2-u^3/3)\le r/6$ for $u\in [0,1]$.

\paragraph{Spatial discretization and admissible set.}
Let $u=(u_1,\dots,u_N)^\top\in\R^N$ denote the spatially discretized state. The admissible set becomes the box
\begin{equation}\label{eq:KN-fisher}
\cK_N=[0,1]^N,
\end{equation}
and the interior set on which the latent construction operates is $(0,1)^N$.

\paragraph{Latent map.}
With $m=0$, $M=1$ the sigmoid parametrization of Definition~\ref{def:latent-map} simplifies to
\begin{equation}\label{eq:g-fisher}
g(z)_i=\sigma(z_i),\qquad g^{-1}(u)_i=\log\frac{u_i}{1-u_i},\qquad i=1,\dots,N.
\end{equation}

\paragraph{Latent energy and induced discrete energy.}
The reaction--diffusion ansatz of Definition~\ref{def:rd-energy} takes the form
\begin{equation}\label{eq:psi-fisher}
\Psi_{\theta}(z)=\sum_{i=1}^N V_{\theta}(z_i)+\frac12\sum_{i,j=1}^N a_{ij}^{(\theta)}(z_i-z_j)^2,
\end{equation}
where $V_{\theta}$ and $a_{ij}^{(\theta)}$ are parametrized as in Definition~\ref{def:neural-param}. The induced discrete energy on $(0,1)^N$ is
\begin{equation}\label{eq:Etheta-fisher}
\cE_{\theta}(u)=\Psi_{\theta}(g^{-1}(u)).
\end{equation}

\paragraph{Architecture guarantees.}
By Proposition~\ref{prop:structure}, assuming the latent ODE generates a global flow (which holds under the conditions of Proposition~\ref{prop:global-flow}), the learned family $\PhiT(\cdot,\tau)$ satisfies admissibility ($\PhiT(u,\tau)\in[0,1]^N$) and induced-energy dissipation $\cE_{\theta}(\PhiT(u,\tau))\le\cE_{\theta}(u)$ exactly on $(0,1)^N$.

\paragraph{Matching condition.}
Let $\Delta_h$ denote the discrete Laplacian corresponding to the spatial discretization and $\odot$ denote componentwise multiplication. The generator matching condition \eqref{eq:matching} then reads
\begin{equation}\label{eq:matching-fisher}
\mathrm{D}g(g^{-1}(u))\,K_{\theta}(g^{-1}(u))\,\nabla\Psi_{\theta}(g^{-1}(u)) = -\bigl(\nu\Delta_h u + r\,u\odot(1-u)\bigr),\qquad u\in(0,1)^N.
\end{equation}
This identity is a target to be realized by the neural parametrization of $\Psi_{\theta}$ and $K_{\theta}$; it is not automatically satisfied by the architecture.

\subsection{Allen--Cahn equation}\label{subsec:allen-cahn}
As a second benchmark, consider the Allen--Cahn equation
\begin{equation}\label{eq:allen-cahn-bench}
\partial_t u = \varepsilon^2 \Delta u + u(1-u^2), \qquad -1\le u\le 1,
\end{equation}
on $\Omega\subset \R^d$ with homogeneous Neumann or periodic boundary conditions. The admissible set is $\cK=\{u:-1\le u\le 1\}$, and the associated energy is
\begin{equation}\label{eq:allen-cahn-energy}
\cE(u)=\int_\Omega\Bigl(\frac{\varepsilon^2}{2}|\nabla u|^2 + \frac{(1-u^2)^2}{4}\Bigr)\,dx,
\end{equation}
which satisfies $\cE(\cS_t u)\le \cE(u)$ for all $t\ge 0$. The admissible set becomes $\cK_N=[-1,1]^N$ with interior $(-1,1)^N$. The latent map, energy ansatz, and matching condition follow the same pattern as for Fisher--KPP with $m=-1$, $M=1$, and the generator matching target
\begin{equation}\label{eq:matching-allen-cahn}
\mathrm{D}g(g^{-1}(u))\,K_{\theta}(g^{-1}(u))\,\nabla\Psi_{\theta}(g^{-1}(u)) = -\bigl(\varepsilon^2\Delta_h u + u\odot(1-u\odot u)\bigr).
\end{equation}

\subsection{Experimental setup}\label{subsec:exp-setup}

\paragraph{Architecture details.}
Table~\ref{tab:architecture} summarizes the network architectures used in all experiments. The latent learner uses the parametrization of Definitions~\ref{def:rd-energy}--\ref{def:neural-param} with a local interaction mask (each site connected to its $2r_K$ nearest neighbors, periodic padding for $K_{\theta}$, periodic interaction graph for $a_{ij}^{(\theta)}$). The activation function $\rho$ is $\softplus$ throughout, which is $C^{\infty}$ and compatible with the regularity requirements of Propositions~\ref{prop:global-flow}--\ref{prop:lipschitz}. The ResNet baseline uses a standard residual network with skip connections, trained with a pointwise MSE loss.

\begin{table}[htbp]
\centering
\caption{Architecture details for the latent semigroup learner, FNO, and ResNet baseline.}
\label{tab:architecture}
\small
\begin{tabular}{@{}lccc@{}}
\hline
Component & Latent learner & FNO & ResNet baseline \\
\hline
Potential $V_{\theta}$ & $L_V{=}2$, widths $(64,64)$, $\beta_V{=}0$${}^{\dagger}$ & --- & --- \\
Interaction rank $r_A$ & $8$ (nearest-neighbor, radius~$2$) & --- & --- \\
Mobility $k_{\theta,i}$ & $L_K{=}2$, widths $(64,64)$, $r_K{=}3$ & --- & --- \\
Fourier layers & --- & $4$ layers, $16$ modes, width~$16$ & --- \\
Residual blocks & --- & --- & $4$ blocks, $32$ ch. \\
Total parameters & $\sim$9.6\,K & $\sim$17.5\,K & $\sim$41.6\,K \\
\hline
Optimizer & AdamW, wd=$10^{-5}$ & AdamW, wd=$10^{-5}$ & AdamW, wd=$10^{-5}$ \\
Grad clip & $1.0$ & $1.0$ & $1.0$ \\
Learning rate & $10^{-3}$ (cosine decay) & $10^{-3}$ (cosine decay) & $10^{-3}$ (cosine decay) \\
Batch size & $64$ & $64$ & $64$ \\
Training epochs & $100$ & $100$ & $100$ \\
\hline
\end{tabular}
\end{table}
${}^{\dagger}$Archived experiments use $\beta_V=0$ and are not covered by Proposition~\ref{prop:global-flow}; finite-trajectory stability is only an empirical diagnostic (Remark~\ref{rem:betaV-zero}).

\paragraph{Training data.}
For each benchmark, training data consist of triples $(u_0,\tau,\cS_{\tau}u_0)$ with $1000$ training samples and $50$ hold-out test samples.
The initial conditions $u_0$ are drawn from PDE-specific distributions within the interior of the admissible set: for Fisher--KPP, $u_0\sim\mathrm{Uniform}[0.3,\,0.7]^N$; for Allen--Cahn, $u_0=\tanh((x-0.5)/(\sqrt{2}\,\varepsilon))+\xi$ with $\xi\sim 0.05\,\cN(0,I_N)$; for Burgers, $u_0=\sin(2\pi x/L)\cdot(0.5+0.5\,\eta)+\xi$ with $\eta\sim\mathrm{Uniform}[0,1]$ and $\xi\sim 0.1\,\cN(0,I_N)$.
The time step is $\tau=\Delta t=0.1$ for Fisher--KPP and Allen--Cahn, and $\tau=\Delta t=0.05$ for the stiffer Burgers benchmark.
The reference solution $\cS_{\tau}u_0$ is computed by a high-resolution numerical integrator with a finer internal time step $\Delta t_{\mathrm{ref}}$: central finite differences with RK4 ($\Delta t_{\mathrm{ref}}=0.005$) for Fisher--KPP, upwind finite differences with RK4 ($\Delta t_{\mathrm{ref}}=0.001$) for Burgers, and a semi-implicit IMEX scheme with spectral treatment of the diffusion term ($\Delta t_{\mathrm{ref}}=0.0005$) for Allen--Cahn.

\paragraph{Computational cost.}
All experiments were conducted on a single NVIDIA GeForce RTX 5080 GPU. The latent learner requires $M_{\mathrm{sub}}=30$ RK4 substeps per time step, corresponding to 120 function evaluations per step (vs.\ 1 for FNO and ResNet). Despite this overhead, the latent learner's total training time is comparable to the baselines due to its smaller parameter count (9.6\,K vs.\ 17.5\,K for FNO and 41.6\,K for ResNet). The per-epoch training times are approximately $45$\,s for the latent learner, $30$\,s for FNO, and $35$\,s for ResNet on the Fisher--KPP benchmark ($N=64$, batch size 64).

\paragraph{Reproducibility.}
The archived runs did not fix every source of randomness. The revised runners seed Python, NumPy, PyTorch, and CUDA from one recorded seed and optionally request deterministic PyTorch algorithms. Linear latent weights retain the initialization $\mathcal{N}(0,0.001)$ with zero biases. Validation-trajectory variation is not a substitute for variation across independently trained seeds; revised comparative claims should therefore report multiple training seeds or explicitly identify a single-seed result.

\subsection{Results}\label{subsec:results}

\paragraph{Main comparison (Fisher--KPP).}
Table~\ref{tab:benchmark-results} reports the evaluation metrics on $50$ hold-out initial conditions, each rolled out over $20$ time steps $(T_{\max}=2.0)$.

\begin{table}[htbp]
\centering
\caption{Archived Fisher--KPP results ($50$ test samples, mean $\pm$ std). The rollout values were produced before strict reference-time alignment was enforced and are retained for provenance, not as revised comparative evidence.}
\label{tab:benchmark-results}
\small
\begin{tabular}{@{}lccc@{}}
\hline
Metric & Latent semigroup & FNO & ResNet baseline \\
\hline
Parameters & 9.6\,K & 17.5\,K & 41.6\,K \\
Rollout MSE (final) & $(1.49\pm 0.15)\times10^{-4}$ & $(2.10\pm 1.07)\times10^{-3}$ & $(2.86\pm 0.22)\times10^{-2}$ \\[2pt]
Boundary violation & $0.00$\footnotemark[1] & $(4.56\pm 12.2)\times10^{-4}$ & $(7.23\pm 11.52)\times10^{-4}$ \\[2pt]
Numerical SG defect (archived MSE) & $\sim 10^{-14}$ & $(2.19\pm 0.83)\times10^{-5}$ & $(2.05\pm 0.15)\times10^{-4}$ \\[2pt]
Learned-energy monotonicity & $100.0\%$ (archived diagnostic) & Not evaluated & Not evaluated \\
\hline
\end{tabular}
\footnotetext{The boundary violation of zero is guaranteed by the sigmoid map $g\colon \R^N\to(m,M)^N$, which maps all states into the interior of the admissible set by construction. This guarantee holds provided the latent ODE generates a global flow (see Proposition~\ref{prop:global-flow} for sufficient conditions).}
\end{table}

The archived table reported a roughly $190$-fold MSE ratio, but the former evaluator advanced the model by $\tau=0.1$ while indexing the dense Fisher--KPP reference trajectory as though adjacent snapshots were separated by the same interval; they were separated by $0.005$. The ratio must therefore be recomputed. The zero bound-violation observation remains consistent with the bounded decoder on finite numerical trajectories, but the archived $\beta_V=0$ configuration is outside the stated coercive global-flow theorem. Likewise, the $100\%$ value concerns $\cE_\theta$, not the physical energy \eqref{eq:fisher-energy}, and the reported composition residual is numerical.

\paragraph{Convergence under spatial refinement.}
Table~\ref{tab:convergence-N} reports the final-step trajectory MSE for the latent learner on Fisher--KPP as the spatial resolution $N$ varies. The error decreases monotonically with $N$, consistent with the increasing expressiveness of the latent architecture. The successive refinement ratios are approximately $7.2\times$ ($N{=}16{\to}32$), $8.1\times$ ($N{=}32{\to}64$), and $3.9\times$ ($N{=}64{\to}128$). The first two doublings are consistent with an effective convergence rate of approximately third order in $1/N$ (since $\log_2(7.5)\approx 2.9$), while the reduced ratio at $N{=}128$ suggests that the optimizer or expressiveness ceiling begins to limit further improvement. A rigorous convergence-rate analysis is deferred to Problem~\ref{prob:expressivity}.

\begin{table}[htbp]
\centering
\caption{Convergence under spatial refinement for the latent learner on Fisher--KPP ($T_{\max}=2.0$, $\Delta t=0.1$; results are from the main experiment code with $N$ varied in the configuration).}
\label{tab:convergence-N}
\begin{tabular}{cc}
\hline
$N$ & Rollout MSE (final step) \\
\hline
$16$ & $(8.72\pm 1.03)\times10^{-3}$ \\
$32$ & $(1.21\pm 0.18)\times10^{-3}$ \\
$64$ & $(1.49\pm 0.15)\times10^{-4}$ \\
$128$ & $(3.84\pm 0.42)\times10^{-5}$ \\
\hline
\end{tabular}
\end{table}

\begin{remark}[Scope of the convergence study]
The convergence study (Table~\ref{tab:convergence-N}) is presented only for Fisher--KPP. For Allen--Cahn and Burgers, the same latent architecture is used with identical hyperparameters; we defer PDE-specific convergence analysis to future work. A temporal convergence study (varying $\Delta t$ at fixed $N$) is also deferred, as the current experiments fix $\Delta t$ based on the PDE timescale.
\end{remark}

\paragraph{Long-time time-stepping.}
Table~\ref{tab:long-time} reports the trajectory MSE at $T=10.0$ ($100$ time steps) and $T=20.0$ ($200$ time steps) for Fisher--KPP with $N=64$, comparing the latent learner, FNO, and ResNet baseline. The ResNet baseline exhibits numerical blowup beyond $T=10$, with MSE reaching $\sim 10^{16}$ at $T=20$. The FNO maintains finite error but with growing boundary violations ($1.8\%$ at $T=20$). The latent learner is the only model that maintains both bounded error and exact boundary compliance throughout.

\begin{table}[htbp]
\centering
\caption{Long-time rollout comparison on Fisher--KPP ($N=64$, $\Delta t=0.1$). The ResNet baseline exhibits numerical blowup beyond $T=10$. The $T{=}2$ row reproduces the short-time comparison from Table~\ref{tab:benchmark-results} ($50$ test samples) for reference; the $T{=}10$ and $T{=}20$ rows are from a separate long-time rollout experiment ($20$ test samples).}
\label{tab:long-time}
\small
\begin{tabular}{@{}lccc@{}}
\hline
$T_{\max}$ (steps) & Latent semigroup & FNO & ResNet \\
\hline
$2.0$ (20) & $(1.49\pm 0.15)\times10^{-4}$ & $(2.10\pm 1.07)\times10^{-3}$ & $(2.86\pm 0.22)\times10^{-2}$ \\
$10.0$ (100) & $(4.09\pm 0.14)\times10^{-1}$ & $(4.52\pm 0.34)\times10^{-1}$ & $1.56\pm 0.96$ \\
$20.0$ (200) & $(4.10\pm 0.15)\times10^{-1}$ & $(4.66\pm 0.33)\times10^{-1}$ & $(8.60\pm 22.0)\times10^{16}$ \\
\hline
\end{tabular}
\end{table}

\paragraph{Boundary violation growth.}
Table~\ref{tab:boundary-growth} reports the mean boundary violation at selected time horizons. The latent learner maintains zero violation at all times by construction. The FNO and ResNet baseline both exhibit growing violations; the ResNet baseline eventually leaves the admissible set entirely.

\begin{table}[htbp]
\centering
\caption{Boundary violation growth on Fisher--KPP ($N=64$, $\Delta t=0.1$; $20$ test samples from the long-time rollout experiment).}
\label{tab:boundary-growth}
\small
\begin{tabular}{@{}lccc@{}}
\hline
$T$ & Latent semigroup & FNO & ResNet \\
\hline
$0.5$ & $0.00$ & $(8.8\pm 2.1)\times10^{-6}$ & $0.00$ \\
$2.0$ & $0.00$ & $(1.8\pm 0.5)\times10^{-3}$ & $(1.3\pm 0.4)\times10^{-3}$ \\
$10.0$ & $0.00$ & $(1.5\pm 0.3)\times10^{-2}$ & $(6.0\pm 1.1)\times10^{-1}$ \\
$20.0$ & $0.00$ & $(1.8\pm 0.4)\times10^{-2}$ & $(8.7\pm 2.3)\times10^{7}$ \\
\hline
\end{tabular}
\end{table}

\paragraph{Allen--Cahn benchmark.}
The Allen--Cahn rerun protocol uses $\varepsilon=0.1$, $N=64$, $\Delta t=0.1$, and $T_{\max}=2.0$. Complete revised metrics are not available in this snapshot, so Table~\ref{tab:allen-cahn-results} records their status rather than displaying fabricated placeholders.

\begin{table}[htbp]
\centering
\caption{Status of the revised Allen--Cahn evaluation ($\varepsilon=0.1$, $N=64$, $\Delta t=0.1$, $T_{\max}=2.0$).}
\label{tab:allen-cahn-results}
\begin{tabular}{lccc}
\hline
Metric & Latent semigroup & FNO & ResNet baseline \\
\hline
Rollout MSE (final step) & Not reported & Not reported & Not reported \\[2pt]
Boundary violation & Not reported & Not reported & Not reported \\[2pt]
Numerical SG defect & Not reported & Not reported & Not reported \\[2pt]
Learned / physical energy & Not reported & Not evaluated & Not evaluated \\
\hline
\end{tabular}
\end{table}

\paragraph{Discussion.}
The archived tables motivate the revised study but do not yet establish comparative accuracy under the corrected evaluator. The bounded decoder gives a clear finite-trajectory constraint mechanism, while claims about PDE accuracy, physical-energy decay, numerical composition error, or spatial rates require the aligned multi-seed reruns described in the accompanying revision plan.

\subsection{Viscous Burgers equation}\label{subsec:burgers}
As an out-of-class transport stress test, consider the viscous Burgers equation
\begin{equation}\label{eq:burgers-bench}
\partial_t u + u\,\partial_x u = \nu\,\partial_{xx} u, \qquad -1\le u\le 1,
\end{equation}
on $[0,2\pi]$ with periodic boundary conditions and $\nu=0.01$. The admissible set is $\cK_N=[-1,1]^N$, and the energy $\cE(u)=\frac{1}{2}\int u^2\,dx$ satisfies $\cE(\cS_t u)\le \cE(u)$ for all $t\ge 0$, since the convective term $\int u\,u_x\,u\,dx=\frac{1}{3}\int \partial_x(u^3)\,dx=0$ under periodic boundary conditions. This benchmark is moderately stiffer than Fisher--KPP: the smaller viscosity $\nu=0.01$ permits near-shock formation, and the nonlinear convection term $u\,\partial_x u$ introduces genuine transport dynamics absent from the reaction--diffusion benchmarks. The Burgers benchmark uses $\Delta t=0.05$ (versus $\Delta t=0.1$ for Fisher--KPP) to accommodate the smaller viscosity coefficient.

The physical $L^2$ energy dissipates under periodic boundary conditions, but the convective term is not a pure gradient contribution of the implemented form. Consequently, neither the transfer theorem nor the expressivity result is claimed for this benchmark. Existing FNO and ResNet values are retained as archived baselines; the missing latent entries are reported as unavailable.

\begin{table}[htbp]
\centering
\caption{Archived baselines and current status for the out-of-class viscous Burgers stress test ($\nu=0.01$, $N=64$, $\Delta t=0.05$, $T_{\max}=1.0$).}
\label{tab:burgers-results}
\small
\begin{tabular}{@{}lccc@{}}
\hline
Metric & Latent semigroup & FNO & ResNet \\
\hline
Rollout MSE (final) & Not reported & $(5.31\pm 3.96)\times10^{-2}$ & $(8.19\pm 2.51)\times10^{-2}$ \\[2pt]
Boundary violation & Not reported & $(1.54\pm 0.54)\times10^{-1}$ & $(6.56\pm 2.00)\times10^{-2}$ \\[2pt]
Numerical SG defect (archived MSE) & Not reported & $(2.66\pm 1.98)\times10^{-4}$ & $(1.30\pm 0.31)\times10^{-3}$ \\[2pt]
Learned / physical energy & Not reported & Not evaluated & Not evaluated \\
\hline
\end{tabular}
\end{table}

Long-time Burgers rollout and an architecture with an explicit skew transport component are deferred to future work.

\subsection{Ablation study}\label{subsec:ablation}
To isolate the contribution of each architectural component, we evaluate three variants of the latent semigroup learner on the Fisher--KPP benchmark (Table~\ref{tab:ablation}). The \emph{full} model uses the complete architecture of Definitions~\ref{def:rd-energy}--\ref{def:neural-param}. The \emph{no-interaction} variant sets the interaction coefficients $a_{ij}^{(\theta)}=0$ in \eqref{eq:rd-energy}, removing spatial coupling. The \emph{no-mobility} variant replaces the learned diagonal mobility $K_{\theta}(z)$ with the identity matrix, reducing the latent dynamics to a plain gradient flow.

\begin{table}[htbp]
\centering
\caption{Ablation study on Fisher--KPP ($N=64$, $\Delta t=0.1$, $T_{\max}=2.0$). All variants maintain exact boundary compliance by construction (under the global-flow assumption of Proposition~\ref{prop:global-flow}; see the footnote on Table~\ref{tab:benchmark-results} for the guarantee). The full-architecture MSE ($1.89\times10^{-4}$) differs from Table~\ref{tab:benchmark-results} ($1.49\times10^{-4}$) due to different random seeds governing the training--validation split and parameter initialization.}
\label{tab:ablation}
\small
\begin{tabular}{@{}lcc@{}}
\hline
Variant & Rollout MSE & Boundary violation \\
\hline
Full architecture & $(1.89\pm 0.18)\times10^{-4}$ & $0.00$ \\
No interaction ($a_{ij}=0$) & $(1.89\pm 0.18)\times10^{-4}$ & $0.00$ \\
No mobility ($K=I$) & $(6.81\pm 4.77)\times10^{-3}$ & $0.00$ \\
\hline
\end{tabular}
\end{table}

In the archived ablation, removing spatial interaction changed the reported Fisher--KPP error little, whereas replacing the learned mobility produced a roughly $36$-fold error ratio. These observations are hypotheses for the aligned rerun, not causal conclusions: the variants must be trained across matched seeds before attributing the difference to a component. All finite numerical trajectories remained inside the decoder range.

\iffalse % Superseded transfer/expressivity draft retained for review provenance.
\section{Continuous-to-discrete structural transfer}\label{sec:transfer-old}

Having demonstrated the architectural guarantees empirically in Section~\ref{sec:benchmark}, we now turn to the rigorous analytical framework. This section provides the analytical bridge between the PDE-level semigroup and the finite-dimensional latent architecture. We formalize the assumptions on the PDE and its semidiscrete approximation, and then prove a transfer theorem that composes the two approximation layers.

\begin{assumption}[PDE well-posedness and structure]\label{ass:pde-structure}
The PDE \eqref{eq:pde} generates a semigroup $(\cS_t)_{t\ge 0}$ on a Banach space $X$ \eqref{eq:semigroup-family} satisfying:
\begin{enumerate}[label=(\roman*)]
    \item \textbf{Well-posedness:} for every $u_0\in X$, the solution $u(t)=\cS_tu_0$ exists globally in time and depends continuously on the initial data.
    \item \textbf{Invariant admissible set:} there exists a closed set $\cK\subset X$ such that $\cS_t(\cK)\subset \cK$ for all $t\ge 0$ \eqref{eq:maximum-principle}.
    \item \textbf{Lyapunov dissipation:} there exists $\cE:X\to \R\cup\{+\infty\}$ such that $\cE(\cS_tu)\le \cE(u)$ for all $t\ge 0$ and $u\in \cK$ \eqref{eq:dissipation}.
\end{enumerate}
\end{assumption}

\begin{assumption}[Semidiscrete spatial discretization]\label{ass:semidiscrete}
There exists a spatial discretization operator $\cFN:\cK_N\to \R^N$ approximating $\cF$ in the following sense:
\begin{enumerate}[label=(\roman*)]
    \item \textbf{Consistency:} for smooth $u\in \cK$, the spatial discretization error satisfies $\norm{\cFN(u_N)-\cF(u)}\le C h^p$ for a suitable projection $u_N$ of $u$, with mesh size $h$ and order $p\ge 1$.
    \item \textbf{Structural preservation:} the semidiscrete scheme
    \begin{equation}\label{eq:semidiscrete-scheme}
    \partial_t u_N = \cFN(u_N), \qquad u_N(0)=u_0^N,
    \end{equation}
    admits a one-step integrator $S_{\Delta t}^N$ (e.g., forward Euler or another SSP integrator) such that:
    \begin{enumerate}[label=(\alph*)]
        \item \emph{Invariance:} $u_N\in \cK_N \Longrightarrow S_{\Delta t}^N u_N\in \cK_N$;
        \item \emph{Dissipation:} $\cE_N(S_{\Delta t}^N u_N)\le \cE_N(u_N)$ for a suitable discrete energy $\cE_N:\cK_N\to \R$.
    \end{enumerate}
    \item \textbf{Convergence:} the semidiscrete flow satisfies $\norm{S_t^N u_0^N-\cS_tu_0}\le C(N)$ for $0\le t\le T$, with $C(N)\to 0$ as $N\to \infty$.
\end{enumerate}
\end{assumption}

\begin{proposition}[Semidiscrete convergence]\label{prop:semidiscrete-conv}
Under Assumptions~\ref{ass:pde-structure} and \ref{ass:semidiscrete}, for initial data $u_0\in \cK$ and $u_0^N$ a suitable projection or interpolation of $u_0$ onto $X_N$, the semidiscrete solution satisfies
\begin{equation}\label{eq:semidiscrete-conv}
\norm{S_t^N u_0^N - \cS_tu_0}\le C(N)+\varepsilon_{\mathrm{sg}}^N(t), \qquad 0\le t\le T,
\end{equation}
where $C(N)$ is the spatial discretization error and $\varepsilon_{\mathrm{sg}}^N(t)$ is the semigroup defect of the semidiscrete scheme, measuring the deviation from exact semigroup composition.
\end{proposition}

\begin{proof}
We decompose the total error via the triangle inequality:
\begin{equation*}
\norm{S_t^N u_0^N - \cS_tu_0} \le \underbrace{\norm{S_t^N u_0^N - \tilde S_t^N u_0^N}}_{\text{time discretization}} + \underbrace{\norm{\tilde S_t^N u_0^N - \cS_tu_0}}_{\text{spatial discretization}},
\end{equation*}
where $\tilde S_t^N$ denotes the exact flow of the semidiscrete ODE \eqref{eq:semidiscrete-scheme} and $S_t^N$ denotes the discrete-time integrator applied over $n=t/\Delta t$ steps.

\paragraph{Spatial error.} By consistency (Assumption~\ref{ass:semidiscrete}(i)), $\norm{\cFN(u_N)-\cF(u)}\le C h^p$ for smooth $u$. Since the semidiscrete scheme inherits the structural preservation from Assumption~\ref{ass:semidiscrete}(ii), the exact flow $\tilde S_t^N$ is Lipschitz on compact subsets of $\cK_N$ with a constant $L_N$ independent of $h$. A standard Gronwall argument gives $\norm{\tilde S_t^N u_0^N - \cS_tu_0}\le C(N)$, where $C(N)\to 0$ as $N\to\infty$ at rate $O(h^p)$.

\paragraph{Time discretization error.} The one-step integrator $S_{\Delta t}^N$ satisfies $\norm{S_{\Delta t}^N v - \tilde S_{\Delta t}^N v}\le C(\Delta t)^{p+1}$ for $v\in\cK_N$, by standard truncation error analysis under the Lipschitz condition on $\cFN$. Summing over $n=t/\Delta t$ steps via a discrete Gronwall inequality yields $\norm{S_t^N u_0^N - \tilde S_t^N u_0^N}\le \varepsilon_{\mathrm{sg}}^N(t)$, where $\varepsilon_{\mathrm{sg}}^N(t)\le C\,t\,(\Delta t)^p$ for an SSP integrator of order $p$.
\end{proof}

\begin{theorem}[Transfer theorem]\label{thm:transfer}
Let $(\cS_t)_{t\ge 0}$ satisfy Assumption~\ref{ass:pde-structure}, and let $\cFN$ satisfy Assumption~\ref{ass:semidiscrete} with convergence rate $C(N)$. Let $\PhiT$ be the latent architecture of Definition~\ref{def:latent-semigroup}, with the generator matching condition \eqref{eq:matching} satisfied up to error $\varepsilon_{\mathrm{match}}$:
\begin{equation}\label{eq:matching-error}
\norm{\mathrm{D}g(g^{-1}(u))K_{\theta}(g^{-1}(u))\nabla \Psi_{\theta}(g^{-1}(u))+\cFN(u)}\le \varepsilon_{\mathrm{match}}, \qquad u\in (m,M)^N.
\end{equation}
Then, for every $u_0\in (m,M)^N$ and every $0\le n\Delta t\le T$:
\begin{enumerate}[label=(\roman*)]
    \item \textbf{Admissibility:} $\PhiT(u_0,n\Delta t)\in [m,M]^N$ (exact, by Proposition~\ref{prop:structure}).
    \item \textbf{Energy dissipation:} $\cE_{\theta}(\PhiT(u_0,n\Delta t))\le \cE_{\theta}(u_0)$ (exact, by Proposition~\ref{prop:structure}).
    \item \textbf{PDE approximation:}
    \begin{equation}\label{eq:transfer-bound}
    \norm{\PhiT(u_0,n\Delta t)-\cS_{n\Delta t}u_0}\le C(N)+\varepsilon_{\mathrm{sg}}^N(n\Delta t)+e^{L_TT}\,T\,\varepsilon_{\mathrm{match}}+\varepsilon_{\mathrm{app}}(N),
    \end{equation}
    where $\varepsilon_{\mathrm{app}}(N)$ accounts for the coupling between the two approximation levels. Specifically, if $u_0^N$ is the projection of $u_0$ onto $X_N$ and $L_T$ is the Lipschitz constant from \eqref{eq:lipschitz-step}, then $\varepsilon_{\mathrm{app}}(N)\le e^{L_TT}\norm{g(g^{-1}(u_0^N))-u_0^N}+C(N)$, where the first term is the interpolation error through the latent map and the second is the amplification of the semidiscrete error through the latent flow.
\end{enumerate}
\end{theorem}

\begin{proof}
Claims (i) and (ii) are direct consequences of Proposition~\ref{prop:structure}, since $\PhiT$ maps $(m,M)^N$ into $[m,M]^N$ by construction and dissipates $\cE_{\theta}$ along the latent flow.

For claim (iii), we decompose the total error into two independent layers via the triangle inequality:
\begin{equation}\label{eq:transfer-decomp}
\norm{\PhiT(u_0,n\Delta t)-\cS_{n\Delta t}u_0}\le \underbrace{\norm{\PhiT(u_0,n\Delta t)-S_{n\Delta t}^Nu_0^N}}_{\text{(A): latent vs.\ semidiscrete}}+\underbrace{\norm{S_{n\Delta t}^Nu_0^N-\cS_{n\Delta t}u_0}}_{\text{(B): semidiscrete vs.\ PDE}}.
\end{equation}

\paragraph{Bound on (B).} By Proposition~\ref{prop:semidiscrete-conv}, $\norm{S_{n\Delta t}^Nu_0^N-\cS_{n\Delta t}u_0}\le C(N)+\varepsilon_{\mathrm{sg}}^N(n\Delta t)$, where $\varepsilon_{\mathrm{sg}}^N(t)$ is the semigroup defect of the one-step integrator. Under the structural preservation in Assumption~\ref{ass:semidiscrete}(ii), the semidiscrete scheme inherits the SSP property, and for an SSP integrator of order $p$ the defect satisfies $\varepsilon_{\mathrm{sg}}^N(t)\le C\,t\,(\Delta t)^p$. For the experiments in Section~\ref{sec:benchmark}, which use $M_{\mathrm{sub}}=30$ RK4 substeps per time step, this defect is at the ODE integrator accuracy level ($\sim10^{-14}$) and is negligible relative to the other error terms.

\paragraph{Bound on (A).} We use a Gronwall argument. Define $\delta_n:=\norm{\PhiT(u_0,n\Delta t)-S_{n\Delta t}^Nu_0^N}$. For the base case, $\delta_0=\norm{u_0-u_0^N}$ (since $\PhiT(u_0,0)=u_0$ and $S_0^N=\mathrm{Id}$), which we absorb into the initial-data error term $\varepsilon_{\mathrm{app}}(N)$. For the inductive step, both the latent flow $\PhiT$ and the semidiscrete flow $S^N$ are Lipschitz on compact subsets of the interior set (by Proposition~\ref{prop:lipschitz} for $\PhiT$ and by the structural preservation in Assumption~\ref{ass:semidiscrete}(ii) for $S^N$). Adding and subtracting $\PhiT(S_{n\Delta t}^Nu_0^N,\Delta t)$ in the definition of $\delta_{n+1}$ gives
\begin{equation*}
\delta_{n+1} \le \norm{\PhiT(\PhiT(u_0,n\Delta t),\Delta t)-\PhiT(S_{n\Delta t}^Nu_0^N,\Delta t)} + \norm{\PhiT(S_{n\Delta t}^Nu_0^N,\Delta t)-S_{\Delta t}^N(S_{n\Delta t}^Nu_0^N)}.
\end{equation*}
The first term is bounded by $(1+L_T\Delta t)\,\delta_n$ by the Lipschitz condition \eqref{eq:lipschitz-step-verified}. For the second term, writing $v=S_{n\Delta t}^Nu_0^N$ and using the Taylor expansions $\PhiT(v,\Delta t)=v+\Delta t\,f_\theta(v)+O((\Delta t)^2)$ and $S_{\Delta t}^N(v)=v+\Delta t\,\cFN(v)+O((\Delta t)^2)$, the generator mismatch \eqref{eq:matching-error} yields $\norm{f_\theta(v)+\cFN(v)}\le\varepsilon_{\mathrm{match}}$, hence the one-step increments differ by at most $\varepsilon_{\mathrm{match}}\,\Delta t+O((\Delta t)^2)$. For fixed $\Delta t$, the $O((\Delta t)^2)$ term is absorbed into $\varepsilon_{\mathrm{match}}\,\Delta t$, giving
\begin{equation*}
\delta_{n+1}\le (1+L_T\Delta t)\delta_n + \varepsilon_{\mathrm{match}}\,\Delta t,
\end{equation*}
where $L_T$ bounds the joint Lipschitz constant. Applying the discrete Gronwall inequality yields
\begin{equation*}
\delta_n \le (1+L_T\Delta t)^n\,\delta_0 + \varepsilon_{\mathrm{match}}\,\Delta t \sum_{j=0}^{n-1}(1+L_T\Delta t)^j.
\end{equation*}
Since $(1+L_T\Delta t)^j\le e^{L_TT}$ for all $j\le n$ and $\sum_{j=0}^{n-1}1=n$, the $\varepsilon_{\mathrm{match}}$ contribution is bounded by $e^{L_TT}\,n\Delta t\,\varepsilon_{\mathrm{match}}$, plus the initial-data contribution absorbed into $\varepsilon_{\mathrm{app}}(N)$. Since $n\Delta t\le T$, this is bounded independently of~$n$ by $e^{L_TT}\,T\,\varepsilon_{\mathrm{match}}$.

Combining bounds (A) and (B) in \eqref{eq:transfer-decomp} gives \eqref{eq:transfer-bound}, where $\varepsilon_{\mathrm{app}}(N)$ collects the initial-data interpolation terms: $\varepsilon_{\mathrm{app}}(N)\le e^{L_TT}\norm{g(g^{-1}(u_0^N))-u_0^N}+C(N)$, with the first term accounting for the latent-map interpolation error and the second for the semidiscrete-to-PDE error amplified through the Lipschitz flow.

The one-step comparison between the continuous-time latent flow $\PhiT(\cdot,\Delta t)$ and the discrete integrator $S^N_{\Delta t}$ incurs both the generator matching error $\varepsilon_{\mathrm{match}}\cdot\Delta t$ and the time truncation error $O((\Delta t)^{p+1})$ of the semidiscrete integrator. For fixed $\Delta t$, the truncation error contributes an additional $O(T\cdot(\Delta t)^p)$ term to the total bound, which we absorb into $\varepsilon_{\mathrm{app}}(N)$.
\end{proof}

\begin{remark}[Two-layer transfer structure]\label{rem:transfer-interpretation}
Theorem~\ref{thm:transfer} decomposes the full approximation chain into two layers:
\begin{enumerate}[label=(\alph*)]
    \item \emph{Layer~1 (architecture $\to$ semidiscrete):} the latent architecture $\PhiT$ approximates the semidiscrete scheme $S^N$ with error controlled by the generator matching defect $\varepsilon_{\mathrm{match}}$. Admissibility and dissipation are guaranteed exactly at this level by Proposition~\ref{prop:structure}.
    \item \emph{Layer~2 (semidiscrete $\to$ PDE):} the semidiscrete scheme $S^N$ approximates the PDE semigroup $\cS_t$ with error $C(N)$, which is standard numerical analysis under consistency and stability (Assumption~\ref{ass:semidiscrete}).
\end{enumerate}
The theorem thus reduces Problem~\ref{prob:transfer} (Section~\ref{sec:analytical-questions}) to two verifiable tasks:~(a)~consistency of the semidiscrete generator $\cFN$ with respect to the PDE generator $\cF$, and (b)~expressivity of the neural parametrization for the matching condition \eqref{eq:matching-error}. The first is a classical question in numerical PDE; the second is addressed in Section~\ref{sec:expressivity}.
\end{remark}

\section{Semigroup expressivity of the latent architecture}\label{sec:expressivity}

We now address the expressivity question: can the neural parametrization of Section~\ref{sec:neural-param} realize the generator matching condition \eqref{eq:matching} to arbitrary accuracy? The following result shows that, for smooth semidiscrete generators whose vector field admits a gradient-flow factorization, the answer is affirmative.

\begin{proposition}[Universal approximation of generators]\label{prop:universal-approx}
Under the neural parametrization of Section~\ref{sec:neural-param} with activation $\rho=\softplus$ (which is $C^{\infty}$), the following holds. For any smooth semidiscrete generator $\cFN:(m,M)^N\to \R^N$ that admits a factorization $\cFN(u) = -\mathrm{D}g(g^{-1}(u))\cdot A(u)\cdot \nabla \varphi(u)$ with $A$ positive definite and $\varphi$ smooth, or more generally for any $\cFN$ whose $\mathrm{D}g^{-1}\cdot \cFN$ is a gradient field on $(m,M)^N$, for any compact set $U\subset (m,M)^N$, and for any $\varepsilon>0$, there exist parameters $\theta=(V_{\theta},a_{ij}^{(\theta)},K_{\theta})$ such that
\begin{equation}\label{eq:universal-approx}
\norm{\mathrm{D}g(g^{-1}(u))K_{\theta}(g^{-1}(u))\nabla \Psi_{\theta}(g^{-1}(u))+\cFN(u)}\le \varepsilon, \qquad u\in U,
\end{equation}
provided the network widths in Definition~\ref{def:neural-param} are sufficiently large.
\end{proposition}

\begin{proof}[Proof sketch]
The argument proceeds in five steps.
\begin{enumerate}
    \item \textbf{Scalar potential approximation.} The MLP $V_{\theta}$ in \eqref{eq:V-mlp} with $\softplus$ activation is a universal approximator for both a function and its derivative: since $\softplus\in C^{\infty}$ and $\softplus'>0$ everywhere, derivative approximation follows from function approximation for smooth activations \cite{Cybenko1989,Hornik1991Approximation}. Concretely, for any smooth $V:\R\\to\R$ and any compact interval $I\subset\R$, there exist widths and parameters such that $\sup_{\\xi\in I}|V_{\theta}(\\xi)-V(\\xi)|\le \varepsilon'$ and $\sup_{\\xi\in I}|V_{\theta}'(\\xi)-V'(\\xi)|\le \varepsilon'$. See also Pinkus \cite{Pinkus1999Approximation} for the general theory of simultaneous approximation by ridge functions.

    \item \textbf{Interaction coefficients.} The symmetric coupling coefficients $a_{ij}^{(\theta)}=m_{ij}\,\softplus(e_i^\top e_j)$ in \eqref{eq:A-param} can represent any symmetric nonnegative coupling pattern on the prescribed interaction graph, by choosing the embeddings $e_i$ appropriately. For constant coefficients $a^*$ on a connected graph, one may set $e_i=v$ for all $i$ with $\|v\|^2=\softplus^{-1}(a^*)$, yielding a valid rank-1 Gram matrix. For general nonnegative coupling patterns on nearest-neighbor graphs (as used in the experiments), the partial Gram matrix specified on the edges always admits a PSD completion by choosing sufficiently large diagonal entries $\|e_i\|^2$.

    \item \textbf{Mobility approximation.} The diagonal mobility $K_{\theta}(z)=\diag(k_{\theta,1}(z),\dots,k_{\theta,N}(z))$ with $k_{\theta,i}$ given by \eqref{eq:k-param} can approximate any positive diagonal matrix-valued function on compact sets, since each $k_{\theta,i}$ is a softplus-bounded universal approximator of the local stencil. The positivity constraint $k_{\theta,i}>0$ is enforced by the outer softplus in \eqref{eq:k-param}.

    \item \textbf{Composition approximation.} The architecture produces vector fields of the form $\mathrm{D}g\cdot K\cdot\nabla\Psi$, which are gradient flows of $\Psi$ in the latent coordinates with metric $K$, pulled back to observable coordinates by $g$. For a target generator $F^{\star}$ whose pullback $\mathrm{D}g^{-1}\cdot F^{\star}$ admits a gradient structure $K_{\mathrm{target}}\cdot\nabla\varphi_{\mathrm{target}}$, the architecture can approximate both $K_{\mathrm{target}}$ and $\varphi_{\mathrm{target}}$ independently. This covers all gradient-type dissipative generators, including the reaction--diffusion benchmarks in this paper. More precisely, the composition $F_{\theta}(u):=\mathrm{D}g(g^{-1}(u))K_{\theta}(g^{-1}(u))\nabla \Psi_{\theta}(g^{-1}(u))$ defines a smooth vector field on $(m,M)^N$ for each fixed $\theta$. Fix a target smooth vector field $F^{\star}:U\to\R^N$ of this form. Since $\mathrm{D}g$ is invertible on $U$ (each diagonal entry $\sigma(z_i)(1-\sigma(z_i))>0$), the matching condition $F_{\theta}(u)=F^{\star}(u)$ is equivalent to
    \begin{equation*}
    K_{\theta}(g^{-1}(u))\nabla \Psi_{\theta}(g^{-1}(u)) = \mathrm{D}g(g^{-1}(u))^{-1}F^{\star}(u), \qquad u\in U.
    \end{equation*}
    The right-hand side is a smooth function of $u$, hence of $z=g^{-1}(u)$ on the compact set $g^{-1}(U)$. By steps~1--3, we can approximate $V'$ and each $a_{ij}$ to arbitrary accuracy, hence $\nabla\Psi_{\theta}(z)$ is approximated to arbitrary accuracy on $g^{-1}(U)$. Likewise, the softplus-bounded mobility $K_{\theta}$ can approximate any positive diagonal target. Since the composition is a product of smooth functions that are individually approximated, the full product $K_{\theta}\nabla\Psi_{\theta}\circ g^{-1}$ is dense in the class of smooth vector fields on $U$ that factor through this product structure. For a generic smooth target $F^{\star}$, this constitutes an approximation result for the restricted class of vector fields realizable by the architecture, not a fully general universality claim.

    \item \textbf{Matching condition.} By steps~1--4, for any smooth semidiscrete generator $\cFN$ in the realizable class (i.e., those whose $\mathrm{D}g^{-1}\cdot \cFN$ is a gradient field), there exist parameters $\theta$ such that $F_{\theta}$ approximates $-\cFN$ uniformly on $U$, which is precisely \eqref{eq:universal-approx}.
\end{enumerate}
\end{proof}

\begin{remark}[Scope of the universal approximation result]\label{rem:ua-scope}
Proposition~\ref{prop:universal-approx} establishes expressivity for the class of \emph{gradient-type} semidiscrete generators---those whose pullback $\mathrm{D}g^{-1}\cdot\cFN$ is a gradient field on $(m,M)^N$. This class includes all reaction--diffusion and viscous Burgers discretizations used in the experiments of Section~\ref{sec:benchmark}, whose spatial discretizations are gradient flows with respect to the corresponding discrete energy. For general (non-gradient) generators---for example, those arising from advection-dominated or hyperbolic PDEs---the expressivity of the architecture remains open (Problem~\ref{prob:expressivity}). The proof sketch above approximates the scalar potential and mobility independently; a formal density statement for the full compositional class $F_\theta = \mathrm{D}g\cdot K_\theta\cdot\nabla\Psi_\theta$ in the space of arbitrary smooth vector fields has not been established.
\end{remark}

\begin{corollary}[Resolution of Problem~\ref{prob:expressivity}]\label{cor:resolve-prob1}
Under the assumptions of Theorem~\ref{thm:transfer}, if the neural parametrization of Section~\ref{sec:neural-param} is sufficiently expressive (as guaranteed by Proposition~\ref{prop:universal-approx} for gradient-type smooth semidiscrete generators), then the latent architecture $\PhiT$ approximates the PDE semigroup $\cS_t$ uniformly on compact time intervals $[0,T]$ and on compact subsets of the interior state set $(m,M)^N$:
\begin{equation}\label{eq:cor-prob1-bound}
\sup_{u_0\in U}\sup_{0\le n\Delta t\le T}\norm{\PhiT(u_0,n\Delta t)-\cS_{n\Delta t}u_0}\le C(N)+\varepsilon_{\mathrm{sg}}^N(n\Delta t)+e^{L_TT}\,T\,\varepsilon_{\mathrm{match}}+\varepsilon_{\mathrm{app}}(N),
\end{equation}
with $\varepsilon_{\mathrm{match}}\to 0$ as the network widths increase, for generators within the class covered by Proposition~\ref{prop:universal-approx} (smooth gradient-type generators whose $\mathrm{D}g^{-1}\cdot \cFN$ is a gradient field on $(m,M)^N$).
\end{corollary}

\fi
\section{Boundary-admissible neural semiflows}\label{sec:boundary-semiflow}

The bounded latent gradient-flow construction of Sections~\ref{sec:latent-architecture}--\ref{sec:neural-param} and the boundary-family construction studied in the Wave~2 experiments enforce different structures.  The former enforces an interior pointwise state range and a learned-energy law, whereas the latter enforces a prescribed discrete spatial boundary relation.  This section gives the finite-dimensional theory for the boundary-family construction and then identifies the additional PDE-dependent hypotheses needed for an accuracy statement.  In particular, no result below treats a spatial boundary condition as a consequence of the abstract semigroup law.

\subsection{Temporal object and scope}

For an autonomous initial-boundary-value problem, uniqueness of forward solutions gives a one-parameter solution semiflow $(S_t)_{t\geq 0}$ satisfying $S_{t+s}=S_t\circ S_s$.  By contrast, an explicitly time-dependent equation
\begin{equation}\label{eq:nonautonomous-pde}
\partial_t u=\mathcal F(t,u)
\end{equation}
generally defines a two-parameter evolution family $U(t,s)$ with
\begin{equation}\label{eq:evolution-family-law}
U(t,r)U(r,s)=U(t,s),\qquad s\leq r\leq t,
\end{equation}
not a one-parameter semigroup on the unaugmented state.  A clock variable can make such a system autonomous on an enlarged state space, but the resulting state and comparison problem are different.  Thus autonomy is a structural hypothesis, rather than a property shared by all time-dependent PDEs.

\subsection{Affine boundary geometry on a fixed grid}

Fix $N\geq3$, a grid spacing $\Delta x>0$, and the Euclidean grid space $X_h=\R^N$ with coordinates $u_0,\ldots,u_{N-1}$.  Each boundary family implemented in \texttt{BoundaryFamilySpec} can be written as
\begin{equation}\label{eq:affine-boundary-set}
M_{B,h}:=\{u\in X_h:C_{B,h}u=b_{B,h}\},\qquad \mathrm{T}_uM_{B,h}=\ker C_{B,h}\quad (u\in M_{B,h}).
\end{equation}
For nonhomogeneous Dirichlet data, the two equations are
\begin{equation}\label{eq:dirichlet-discrete-boundary}
C_{D,h}u=(u_0,u_{N-1})^\top,\qquad b_{D,h}=(g_L,g_R)^\top.
\end{equation}
For the homogeneous Neumann family, the code uses first-order outward one-sided differences,
\begin{equation}\label{eq:neumann-discrete-boundary}
C_{N,h}u=\left(-\frac{u_1-u_0}{\Delta x},\frac{u_{N-1}-u_{N-2}}{\Delta x}\right)^\top,\qquad b_{N,h}=0.
\end{equation}
For Robin data with common coefficients $\alpha,\beta$ and endpoint values $\gamma_L,\gamma_R$, it uses
\begin{equation}\label{eq:robin-discrete-boundary}
C_{R,h}u=\left(\left(\alpha+\frac{\beta}{\Delta x}\right)u_0-\frac{\beta}{\Delta x}u_1,\left(\alpha+\frac{\beta}{\Delta x}\right)u_{N-1}-\frac{\beta}{\Delta x}u_{N-2}\right)^\top,
\end{equation}
with $b_{R,h}=(\gamma_L,\gamma_R)^\top$ and the nondegeneracy condition
\begin{equation}\label{eq:robin-nondegenerate}
\alpha+\frac{\beta}{\Delta x}\neq0.
\end{equation}
Equations~\eqref{eq:neumann-discrete-boundary} and \eqref{eq:robin-discrete-boundary} are definitions of the tested discrete endpoint relations.  A claim that they converge to a continuous Neumann or Robin trace requires a separate boundary-consistency estimate as $\Delta x\to0$.
The algebraic statements below require only $N\geq3$; the frozen Wave~2 experiment configuration additionally requires an odd $N\geq5$.

Let $\Pi_{B,h}:X_h\to X_h$ denote the affine state reconstruction implemented by \texttt{project\_state}, and let $Q_{B,h}:X_h\to X_h$ denote the homogeneous tangent reconstruction implemented by \texttt{project\_tangent}.  For Robin data, set
\begin{equation}\label{eq:robin-reconstruction-constants}
q_h:=\frac{\beta}{\Delta x},\qquad d_h:=\alpha+q_h,\qquad \rho_h:=\frac{q_h}{d_h}.
\end{equation}
The endpoint formulas are
\begin{equation}\label{eq:boundary-retraction-formulas}
\begin{aligned}
&\text{Dirichlet:} &&(\Pi_{D,h}u)_0=g_L,\qquad (\Pi_{D,h}u)_{N-1}=g_R,\\
&&&(Q_{D,h}v)_0=(Q_{D,h}v)_{N-1}=0,\\
&\text{Neumann:} &&(\Pi_{N,h}u)_0=u_1,\qquad (\Pi_{N,h}u)_{N-1}=u_{N-2},\\
&&&(Q_{N,h}v)_0=v_1,\qquad (Q_{N,h}v)_{N-1}=v_{N-2},\\
&\text{Robin:} &&(\Pi_{R,h}u)_0=\frac{\gamma_L+q_hu_1}{d_h},\qquad (\Pi_{R,h}u)_{N-1}=\frac{\gamma_R+q_hu_{N-2}}{d_h},\\
&&&(Q_{R,h}v)_0=\rho_hv_1,\qquad (Q_{R,h}v)_{N-1}=\rho_hv_{N-2}.
\end{aligned}
\end{equation}
All unlisted interior entries are unchanged.

\begin{proposition}[Affine boundary and tangent retractions]\label{prop:boundary-retractions}
For each family in \eqref{eq:dirichlet-discrete-boundary}--\eqref{eq:robin-discrete-boundary}, subject to \eqref{eq:robin-nondegenerate}, the maps in \eqref{eq:boundary-retraction-formulas} satisfy
\begin{equation}\label{eq:boundary-retraction-identities}
C_{B,h}\Pi_{B,h}u=b_{B,h},\qquad C_{B,h}Q_{B,h}v=0,
\end{equation}
for all $u,v\in X_h$.  Moreover,
\begin{equation}\label{eq:boundary-retraction-idempotence}
\Pi_{B,h}^2=\Pi_{B,h},\qquad Q_{B,h}^2=Q_{B,h},\qquad \Pi_{B,h}u-\Pi_{B,h}v=Q_{B,h}(u-v).
\end{equation}
Consequently, $\operatorname{Range}(\Pi_{B,h})=M_{B,h}$ and $\operatorname{Range}(Q_{B,h})=\ker C_{B,h}$.
\end{proposition}

\begin{proof}
For Dirichlet data, substitution of the reconstructed endpoint values gives $C_{D,h}\Pi_{D,h}u=b_{D,h}$, while zero tangent endpoints give $C_{D,h}Q_{D,h}v=0$.  For Neumann data, the reconstructed endpoint and its adjacent interior value coincide, so both one-sided differences in \eqref{eq:neumann-discrete-boundary} vanish.  For Robin data, the left endpoint obeys
\begin{equation*}
d_h(\Pi_{R,h}u)_0-q_hu_1=\gamma_L,
\end{equation*}
and the right endpoint obeys the analogous identity; replacing $\gamma_L,\gamma_R$ by zero proves the tangent relation.  A second reconstruction changes neither an admissible endpoint nor an already homogeneous tangent endpoint, proving idempotence.  Finally, subtracting two affine endpoint formulas cancels the boundary data and leaves exactly the corresponding homogeneous formulas.  The range statements follow from \eqref{eq:boundary-retraction-identities}--\eqref{eq:boundary-retraction-idempotence}.
\end{proof}

\begin{remark}[Retraction, not orthogonal projection]\label{rem:boundary-not-orthogonal}
Neither $\Pi_{B,h}$ nor $Q_{B,h}$ is asserted to be orthogonal in the Euclidean or a mesh-weighted inner product.  The properties used below are affine reconstruction, idempotence, and tangency.  Calling these maps orthogonal projections would require an additional inner-product calculation that the implementation neither needs nor supplies.
\end{remark}

\subsection{Global boundary-preserving neural semiflow}

Fix any positive scalar map-duration argument.  In hard autonomous mode, its value is ignored after API validation.  Combine the zero-padded three-point patch-extraction map, the constant temporal feature, and the shared pointwise Tanh multilayer perceptron into a raw vector field $G_{\theta,h}:X_h\to X_h$.  The implemented boundary-compatible vector field is
\begin{equation}\label{eq:hard-autonomous-generator}
F_{\theta,h}(u):=Q_{B,h}G_{\theta,h}(\Pi_{B,h}u).
\end{equation}
The duration requested from the integrator does not enter $G_{\theta,h}$ in autonomous mode.

\begin{theorem}[Boundary-admissible global neural semiflow]\label{thm:boundary-global-semiflow}
Fix the grid, a boundary family satisfying the conditions above, and finite network parameters $\theta$.  Suppose $G_{\theta,h}$ is the finite Tanh network implemented by \texttt{BoundaryFamilyFlow} in hard autonomous mode, namely \texttt{enforcement\_mode=hard} and \texttt{temporal\_mode=autonomous}.  Then:
\begin{enumerate}[label=(\roman*)]
    \item $F_{\theta,h}$ is globally Lipschitz on $X_h$;
    \item for every $u_0\in X_h$, the ODE
    \begin{equation}\label{eq:boundary-neural-ode}
    \dot u=F_{\theta,h}(u),\qquad u(0)=u_0,
    \end{equation}
    has a unique global forward solution;
    \item $M_{B,h}$ is invariant: if $u_0\in M_{B,h}$, then $u(t)\in M_{B,h}$ for all $t\geq0$;
    \item the exact solution maps $\Phi_t^{\theta,h}u_0:=u(t;u_0)$ restrict to a forward semigroup on $M_{B,h}$:
    \begin{equation}\label{eq:boundary-exact-semigroup}
    \Phi_0^{\theta,h}=\operatorname{Id},\qquad \Phi_{t+s}^{\theta,h}=\Phi_t^{\theta,h}\circ\Phi_s^{\theta,h},\qquad s,t\geq0.
    \end{equation}
\end{enumerate}
\end{theorem}

\begin{proof}
Write $\Pi_{B,h}u=Q_{B,h}u+r_{B,h}$, as follows from \eqref{eq:boundary-retraction-idempotence}.  Zero-padded patch extraction is a finite-dimensional linear map, the added autonomous feature is constant, and Tanh is globally $1$-Lipschitz.  Let $U_h$ be the three-point patch operator, let $W_1,W_2,W_3$ be the combined linear maps induced by the three shared MLP layers on the stacked grid features, and take all operator norms in the corresponding Euclidean product spaces.  One admissible bound is
\begin{equation}\label{eq:boundary-generator-lipschitz}
\operatorname{Lip}(F_{\theta,h})\leq \lVert Q_{B,h}\rVert^2\lVert W_3\rVert\lVert W_2\rVert\lVert W_1\rVert\lVert U_h\rVert=:L_{\theta,h}<\infty.
\end{equation}
Biases and the constant temporal feature cancel when two states are subtracted.  This proves global Lipschitz continuity.

Picard--Lindel\"of gives a unique local solution.  Global Lipschitz continuity also gives linear growth,
\begin{equation*}
\lVert F_{\theta,h}(u)\rVert\leq \lVert F_{\theta,h}(0)\rVert+L_{\theta,h}\lVert u\rVert,
\end{equation*}
so Gronwall's inequality bounds the solution on every finite time interval.  Hence no finite-time escape is possible, and the solution extends to every $t\geq0$.

By \eqref{eq:boundary-retraction-identities} and \eqref{eq:hard-autonomous-generator},
\begin{equation*}
C_{B,h}F_{\theta,h}(u)=0
\end{equation*}
for every $u$.  Along a solution beginning in $M_{B,h}$,
\begin{equation*}
\frac{d}{dt}\bigl(C_{B,h}u(t)-b_{B,h}\bigr)=C_{B,h}F_{\theta,h}(u(t))=0,
\end{equation*}
and the initial residual is zero.  This proves invariance.  Finally, the time-shifted curve $u(t+s;u_0)$ and the curve $u(t;u(s;u_0))$ solve the same autonomous ODE with the same initial value at time zero.  Uniqueness yields \eqref{eq:boundary-exact-semigroup}.
\end{proof}

\begin{remark}[Fixed-grid and fixed-parameter guarantee]\label{rem:fixed-grid-guarantee}
The constant in \eqref{eq:boundary-generator-lipschitz} may depend on the trained weights, grid, boundary reconstruction, and chosen norm.  Thus \Cref{thm:boundary-global-semiflow} is a global theorem for each fixed finite-dimensional model, not a stability estimate uniform in mesh refinement or training.  It also proves neither a maximum principle nor physical-energy decay for the Wave~2 Tanh generator; those are separate PDE-specific structures.
\end{remark}

\begin{remark}[Query-time conditioning]\label{rem:query-time-family}
In query-time mode, the requested duration $\tau$ enters the raw field, producing a family $F_{\theta,h}^{[\tau]}$.  For each fixed scalar $\tau$, standard ODE theory may define a flow generated by that fixed vector field, but the requested maps for different durations are not guaranteed to be restrictions of one common generator.  If durations vary across a batch, the implementation evaluates distinct conditioned fields sample by sample; the batched operation is likewise not one common semiflow.  Hence the architecture does not imply $\Phi_{t+s}=\Phi_t\circ\Phi_s$ across queries.  Moreover, a requested duration is not an absolute clock or starting time, so this conditioning does not by itself implement the evolution family in \eqref{eq:evolution-family-law}.
\end{remark}

\subsection{Boundary preservation and composition error of explicit Runge--Kutta maps}

\begin{proposition}[Affine boundary preservation by explicit Runge--Kutta methods]\label{prop:explicit-rk-boundary}
Let $F:X_h\to\ker C_{B,h}$ and let $u_n\in M_{B,h}$.  For any explicit Runge--Kutta method with stages
\begin{equation}\label{eq:explicit-rk-stages}
Y_i=u_n+\delta\sum_{j<i}a_{ij}K_j,\qquad K_i=F(Y_i),
\end{equation}
and update
\begin{equation}\label{eq:explicit-rk-update}
u_{n+1}=u_n+\delta\sum_i b_iK_i,
\end{equation}
every stage and the update belong to $M_{B,h}$ in exact arithmetic.  In particular, the RK4 stages and output used by \texttt{BoundaryFamilyFlow} preserve all three discrete boundary families.
\end{proposition}

\begin{proof}
Proceed by induction over the stages.  Since $C_{B,h}u_n=b_{B,h}$ and $C_{B,h}K_j=0$, \eqref{eq:explicit-rk-stages} gives $C_{B,h}Y_i=b_{B,h}$ whenever the preceding stages are defined.  The same calculation in \eqref{eq:explicit-rk-update} gives $C_{B,h}u_{n+1}=b_{B,h}$.
\end{proof}

\begin{remark}[Role of the repeated reconstruction]\label{rem:rk-reconstruction-role}
On admissible data and in exact arithmetic, the additional \texttt{project\_state} calls at each RK4 stage are identities by \Cref{prop:boundary-retractions,prop:explicit-rk-boundary}.  They therefore do not alter the classical RK4 update on $M_{B,h}$.  In floating-point computation they remove accumulated endpoint residuals, and for an inadmissible incoming state the initial reconstruction places the computation on $M_{B,h}$.  Boundary preservation does not imply that the resulting finite-step map is an exact semigroup.
\end{remark}

For a partition $\pi=(\delta_0,\ldots,\delta_{m-1})$ with $\delta_j>0$ and $\sum_j\delta_j=T$, let $\Psi_{\pi}^{\theta,h}$ denote the composition of RK4 steps for the autonomous field \eqref{eq:hard-autonomous-generator}, and let $|\pi|=\max_j\delta_j$.

\begin{theorem}[RK4 error and nonuniform composition defect]\label{thm:rk4-composition-defect}
Suppose $F_{\theta,h}$ is $C^5$ on a neighborhood of the relevant states up to time $T$, its relevant derivatives are bounded there, and the exact solution, all numerical step values, and all RK stage values remain in that neighborhood for $|\pi|\leq\delta_0$.  Let $R_\delta$ denote one RK4 step and assume the one-step stability estimate
\begin{equation}\label{eq:rk4-one-step-stability}
\lVert R_\delta(v)-R_\delta(w)\rVert\leq(1+C_{\mathrm{stab}}\delta)\lVert v-w\rVert
\end{equation}
holds there for $0<\delta\leq\delta_0$.  Then there is a constant $C_T$, independent of the particular partition, such that
\begin{equation}\label{eq:rk4-nonuniform-global-error}
\lVert\Psi_{\pi}^{\theta,h}(u)-\Phi_T^{\theta,h}(u)\rVert\leq C_T\sum_{j=0}^{m-1}\delta_j^5\leq C_TT|\pi|^4.
\end{equation}
If $\pi$ and $\rho$ are two direct or composed partitions of the same total time $T$, then
\begin{equation}\label{eq:rk4-partition-defect}
\lVert\Psi_{\pi}^{\theta,h}(u)-\Psi_{\rho}^{\theta,h}(u)\rVert\leq C_TT\bigl(|\pi|^4+|\rho|^4\bigr).
\end{equation}
For the fixed finite Tanh field of \Cref{thm:boundary-global-semiflow}, the required smoothness holds; the estimate remains a finite-time, fixed-model statement.
\end{theorem}

\begin{proof}
Classical RK4 consistency gives a local truncation error bounded by $C\delta_j^5$.  Stability of one step gives an error recursion of the form
\begin{equation*}
e_{j+1}\leq(1+C\delta_j)e_j+C\delta_j^5,
\end{equation*}
where $e_j$ is the error at the $j$th partition point.  Discrete Gronwall yields the first inequality in \eqref{eq:rk4-nonuniform-global-error}.  Since $\delta_j^5\leq |\pi|^4\delta_j$ and $\sum_j\delta_j=T$, the second inequality follows.  Both numerical paths in \eqref{eq:rk4-partition-defect} approximate the same exact state $\Phi_T^{\theta,h}(u)$, so the triangle inequality and \eqref{eq:rk4-nonuniform-global-error} prove the result.  By \Cref{rem:rk-reconstruction-role}, the affine stage reconstructions are identities on the exact-arithmetic stages and do not reduce the order in this setting.
\end{proof}

\begin{remark}[Absolute and relative defects]\label{rem:relative-defect-denominator}
Equation~\eqref{eq:rk4-partition-defect} controls an absolute defect.  A relative defect inherits the same order only when its normalizing denominator is bounded below by a specified $\kappa>0$.  If the metric instead divides by $\max\{d,\delta_{\mathrm{floor}}\}$, the immediate bound is \eqref{eq:rk4-partition-defect} divided by $\delta_{\mathrm{floor}}$.  A small relative number without its denominator is not an independent theorem.  For query-time fields, solver refinement removes the RK error but may leave a nonzero mismatch between flows generated by different $F_{\theta,h}^{[\tau]}$.
\end{remark}

\subsection{Generator consistency, stability, and PDE accuracy}

Exact composition can hold for the flow of an incorrect generator.  The missing accuracy bridge is therefore a consistency estimate relative to a target method-of-lines generator together with a stability estimate.

\begin{theorem}[One-sided stability--consistency estimate]\label{thm:generator-stability}
Let $(X_h,\langle\cdot,\cdot\rangle_h)$ be a finite-dimensional Hilbert space.  Let $A_h$ be a target semidiscrete PDE generator and $F_{\theta,h}$ a learned autonomous generator.  Suppose the two solutions $x'=A_h(x)$ and $y'=F_{\theta,h}(y)$ remain in a common set $V_h$ on $[0,T]$, and assume
\begin{equation}\label{eq:one-sided-lipschitz}
\langle F_{\theta,h}(v)-F_{\theta,h}(w),v-w\rangle_h\leq\omega_h\lVert v-w\rVert_h^2,
\end{equation}
for $v,w\in V_h$, together with the generator consistency bound
\begin{equation}\label{eq:generator-consistency}
\sup_{v\in V_h}\lVert F_{\theta,h}(v)-A_h(v)\rVert_h\leq\varepsilon_{\mathrm{gen},h}.
\end{equation}
Define
\begin{equation}\label{eq:chi-omega}
\chi_\omega(t):=\begin{cases}(e^{\omega t}-1)/\omega,&\omega\neq0,\\ t,&\omega=0.\end{cases}
\end{equation}
Then
\begin{equation}\label{eq:stability-consistency-bound}
\lVert y(t)-x(t)\rVert_h\leq e^{\omega_ht}\lVert y(0)-x(0)\rVert_h+\varepsilon_{\mathrm{gen},h}\chi_{\omega_h}(t),\qquad 0\leq t\leq T.
\end{equation}
In particular, if $\omega_h=-\lambda_h<0$ and the initial values agree, then
\begin{equation}\label{eq:contractive-generator-bound}
\lVert y(t)-x(t)\rVert_h\leq\frac{\varepsilon_{\mathrm{gen},h}}{\lambda_h}(1-e^{-\lambda_ht})\leq\frac{\varepsilon_{\mathrm{gen},h}}{\lambda_h}.
\end{equation}
\end{theorem}

\begin{proof}
Set $e=y-x$.  Adding and subtracting $F_{\theta,h}(x)$, and using \eqref{eq:one-sided-lipschitz}--\eqref{eq:generator-consistency}, gives
\begin{equation*}
\frac12\frac{d}{dt}\lVert e\rVert_h^2=\langle F_{\theta,h}(y)-F_{\theta,h}(x),e\rangle_h+\langle F_{\theta,h}(x)-A_h(x),e\rangle_h\leq\omega_h\lVert e\rVert_h^2+\varepsilon_{\mathrm{gen},h}\lVert e\rVert_h.
\end{equation*}
At times when $\lVert e\rVert_h>0$, division by $\lVert e\rVert_h$ yields
\begin{equation*}
\frac{d}{dt}\lVert e\rVert_h\leq\omega_h\lVert e\rVert_h+\varepsilon_{\mathrm{gen},h}.
\end{equation*}
The same inequality holds in the upper-Dini-derivative sense when $e=0$.  Multiplication by $e^{-\omega_ht}$ and integration proves \eqref{eq:stability-consistency-bound}; substituting $\omega_h=-\lambda_h$ proves \eqref{eq:contractive-generator-bound}.
\end{proof}

\begin{corollary}[Layered PDE rollout bound]\label{cor:layered-pde-rollout}
Let $S_t^h$ denote the exact flow of the target semidiscrete system $\dot x=A_h(x)$, and assume it exists through $P_hu_0$ on $[0,T]$.  Let $P_h:X\to X_h$ and $I_h:X_h\to X$ satisfy $\lVert I_hv\rVert_X\leq C_I\lVert v\rVert_h$.  Assume that the target PDE semiflow and exact semidiscrete flow obey
\begin{equation}\label{eq:pde-spatial-error-boundary-section}
\sup_{0\leq t\leq T}\lVert I_hS_t^h(P_hu_0)-S_tu_0\rVert_X\leq\varepsilon_{\mathrm{space}}(h),
\end{equation}
and that a numerical learned state $\widehat y(T)$ satisfies $\lVert\widehat y(T)-y(T)\rVert_h\leq\eta_{\mathrm{time}}(T)$.  Assume the exact learned trajectory $y(t)$ and $x(t)=S_t^h(P_hu_0)$ remain in the common set $V_h$ and satisfy the hypotheses of \Cref{thm:generator-stability}.  Then
\begin{equation}\label{eq:full-layered-error-bound}
\lVert I_h\widehat y(T)-S_Tu_0\rVert_X\leq\varepsilon_{\mathrm{space}}(h)+C_I\left(\eta_{\mathrm{time}}(T)+e^{\omega_hT}\lVert y(0)-P_hu_0\rVert_h+\varepsilon_{\mathrm{gen},h}\chi_{\omega_h}(T)\right).
\end{equation}
\end{corollary}

\begin{proof}
Insert $I_hy(T)$ and $I_hS_T^h(P_hu_0)$, use the triangle inequality and reconstruction stability, apply \eqref{eq:stability-consistency-bound}, and then use \eqref{eq:pde-spatial-error-boundary-section}.
\end{proof}

\begin{remark}[What the current boundary architecture does not guarantee]\label{rem:boundary-current-missing-stability}
The hard autonomous Tanh architecture satisfies \Cref{thm:boundary-global-semiflow}, but its current parametrization does not impose \eqref{eq:one-sided-lipschitz}, \eqref{eq:generator-consistency}, a maximum principle, mass conservation, or decay of a specified physical energy.  Consequently, exact boundary compatibility and a small numerical composition defect do not by themselves make the terms on the right-hand side of \eqref{eq:full-layered-error-bound} small.  This explains how a model can be composition-consistent yet consistently approximate the wrong PDE dynamics.
\end{remark}

\subsection{Why the required structure depends on the PDE}

Autonomy supplies the temporal composition law, but useful invariants and stability constants come from the generator.  The following elementary identities show why different PDE classes require different neural-generator geometries.

\begin{proposition}[Generator geometries for dissipation and conservation]\label{prop:pde-generator-geometries}
Let $E_h\in C^1(X_h)$.
\begin{enumerate}[label=(\roman*)]
    \item If $F_h(u)=-M_h(u)\nabla E_h(u)$ with $M_h(u)$ symmetric positive semidefinite, then
    \begin{equation}\label{eq:l2-gradient-dissipation}
    \frac{d}{dt}E_h(u(t))=-\nabla E_h(u)^\top M_h(u)\nabla E_h(u)\leq0.
    \end{equation}
    \item Let $D_h\mathbf{1}=0$.  If $F_h(u)=-D_h^\top M_h(u)D_h\nabla E_h(u)$ with $M_h(u)$ symmetric positive semidefinite, then both
    \begin{equation}\label{eq:conservative-gradient-identities}
    \frac{d}{dt}\mathbf{1}^\top u(t)=0,\qquad \frac{d}{dt}E_h(u(t))=-\bigl(D_h\nabla E_h(u)\bigr)^\top M_h(u)D_h\nabla E_h(u)\leq0
    \end{equation}
    hold.
    \item If $F_h(u)=(J_h(u)-K_h(u))\nabla E_h(u)$ with $J_h(u)^\top=-J_h(u)$ and $K_h(u)$ symmetric positive semidefinite, then
    \begin{equation}\label{eq:skew-dissipative-identity}
    \frac{d}{dt}E_h(u(t))=-\nabla E_h(u)^\top K_h(u)\nabla E_h(u)\leq0.
    \end{equation}
\end{enumerate}
\end{proposition}

\begin{proof}
For each case, differentiate $E_h(u(t))$ and substitute the stated generator.  In part~(ii),
\begin{equation*}
\mathbf{1}^\top F_h(u)=-(D_h\mathbf{1})^\top M_h(u)D_h\nabla E_h(u)=0.
\end{equation*}
In part~(iii), $z^\top J_hz=0$ for every $z$ because $J_h$ is skew-symmetric.  The remaining quadratic forms have the claimed signs by positive semidefiniteness.
\end{proof}

\Cref{prop:pde-generator-geometries} gives architecture-level algebra once the correct discrete operators have been chosen, but choosing them is PDE-specific.  A reaction--diffusion problem may use an $L^2$ gradient or invariant-region structure; Cahn--Hilliard requires a conservative $H^{-1}$-type operator as in \eqref{eq:conservative-gradient-identities}; transport--diffusion requires a conservative or skew component in addition to diffusion; and a nonautonomous problem requires \eqref{eq:evolution-family-law} or an augmented clock state.  Boundary conditions enter the domains, reconstructions, flux operators, and summation-by-parts identities used in each of these constructions.

\begin{table}[htbp]
\centering
\caption{PDE-dependent theoretical obligations before an experiment can support a transfer claim.  Entries describe the required analysis, not guarantees already supplied by the current Wave~2 model.}
\label{tab:pde-theory-obligations}
\small
\begin{tabular}{@{}>{\raggedright\arraybackslash}p{0.16\textwidth}>{\raggedright\arraybackslash}p{0.16\textwidth}>{\raggedright\arraybackslash}p{0.24\textwidth}>{\raggedright\arraybackslash}p{0.31\textwidth}@{}}
\hline
Problem class & Temporal object & PDE-specific law & Required discrete or neural structure \\
\hline
Fisher--KPP / reaction--diffusion & Forward semiflow & Positivity or invariant interval under compatible data; boundary domain & Boundary-tangent diffusion and a proved invariant-region or comparison mechanism \\
Allen--Cahn & Forward semiflow & Physical free-energy decay & $L^2$ gradient structure matched to the discrete physical energy \\
Cahn--Hilliard & Forward semiflow & Mass conservation and free-energy decay & Conservative $-D_h^\top M_hD_h$ mobility with compatible no-flux or periodic operators \\
Viscous Burgers & Forward semiflow & Conservative transport plus viscous dissipation & Conservative/skew transport component and a dissipative diffusion component \\
Explicitly time-dependent PDE & Evolution family $U(t,s)$ & Start-time and clock dependence & Two-time model or an autonomous model on a correctly augmented state \\
\hline
\end{tabular}
\end{table}

The project therefore uses a PDE theorem card before promoting a new equation to a transfer experiment.  The card must specify the continuous state space and boundary domain, the correct temporal object, well-posedness and stability, the physical invariant or Lyapunov law, the semidiscrete generator and boundary consistency, reconstruction stability, generator matching, time-integration error, and a falsification rule.  The operational template is recorded in \texttt{docs/research/PDE\_THEOREM\_CARD.md}.

\begin{remark}[Claim hierarchy]\label{rem:boundary-claim-hierarchy}
The results of this section have three distinct logical levels.  First, \Cref{prop:boundary-retractions,thm:boundary-global-semiflow,prop:explicit-rk-boundary} are unconditional fixed-grid architecture results for the stated implementation.  Second, \Cref{thm:rk4-composition-defect} is a numerical-analysis result under explicit smoothness and stability hypotheses.  Third, \Cref{thm:generator-stability,cor:layered-pde-rollout} becomes a PDE-accuracy statement only after a PDE-specific method-of-lines, common trajectory set, spatial convergence estimate, and generator error have been supplied.  None of these levels permits the inference ``small composition defect implies small prediction error.''
\end{remark}

\section{Continuous-to-discrete structural transfer}\label{sec:transfer}

This section keeps spatial discretization, exact finite-dimensional flows, and numerical time integration as three distinct objects. Let $P_N:X\to X_N$ be a projection or sampling map and let $I_N:X_N\to X$ be a reconstruction map. We write $\widetilde S_t^N$ for the exact flow generated by
\begin{equation}\label{eq:semidiscrete-scheme-revised}
\dot v=\cFN(v),\qquad v(0)=v_0\in X_N,
\end{equation}
and $Q_h^N$ for a numerical one-step approximation to that flow. Neither symbol is used for both objects.

\begin{assumption}[Exact semidiscrete approximation]\label{ass:semidiscrete}
Fix $T>0$ and a compact set $\cU\subset\cK$. Assume there is a compact common trajectory set $\cV_N\subset(m,M)^N$ that contains $P_N\cU$, $\widetilde S_t^N(P_N\cU)$, and $\PhiT(P_N\cU,t)$ for $0\le t\le T$, and that:
\begin{enumerate}[label=(\roman*)]
    \item $\cFN$ is locally Lipschitz and is $L_N$-Lipschitz on a neighborhood of $\cV_N$;
    \item the reconstruction is stable, $\norm{I_Nv}_X\le C_I\norm{v}_N$;
    \item the exact semidiscrete flow converges after reconstruction:
    \begin{equation}\label{eq:semidiscrete-conv-revised}
    \sup_{u_0\in\cU}\sup_{0\le t\le T}
    \norm{I_N\widetilde S_t^N(P_Nu_0)-\cS_tu_0}_X
    \le \varepsilon_{\mathrm{space}}(N),
    \end{equation}
    where $\varepsilon_{\mathrm{space}}(N)\to0$ as $N\to\infty$.
\end{enumerate}
The PDE-specific projection, reconstruction, norm, invariant neighborhood, and proof or citation for \eqref{eq:semidiscrete-conv-revised} must be supplied when the theorem is applied.
\end{assumption}

On the decoder image $g(\R^N)=(m,M)^N$, define the observable-space generator of the exact latent flow by
\begin{equation}\label{eq:observable-generator}
F_\theta(v):=-\mathrm{D}g(g^{-1}(v))K_\theta(g^{-1}(v))\nabla\Psi_\theta(g^{-1}(v)).
\end{equation}
The use of $g^{-1}$ in \eqref{eq:observable-generator} is specific to the invertible sigmoid parametrization of \Cref{def:latent-map}; it is not available for an arbitrary bounded decoder.

\begin{theorem}[Conditional exact-flow transfer]\label{thm:transfer}
Under Assumption~\ref{ass:semidiscrete}, suppose the exact latent flow exists on $[0,T]$ for the stated initial data and
\begin{equation}\label{eq:matching-error-revised}
\sup_{v\in\cV_N}\norm{F_\theta(v)-\cFN(v)}_N
\le\varepsilon_{\mathrm{match}}.
\end{equation}
Let
\begin{equation*}
\chi_L(t):=
\begin{cases}(e^{Lt}-1)/L,&L>0,\\ t,&L=0.\end{cases}
\end{equation*}
Then, for $u_0\in\cU$ and $0\le t\le T$,
\begin{equation}\label{eq:exact-flow-comparison}
\norm{\PhiT(P_Nu_0,t)-\widetilde S_t^N(P_Nu_0)}_N
\le \varepsilon_{\mathrm{match}}\chi_{L_N}(t),
\end{equation}
and consequently
\begin{equation}\label{eq:transfer-bound}
\norm{I_N\PhiT(P_Nu_0,t)-\cS_tu_0}_X
\le C_I\varepsilon_{\mathrm{match}}\chi_{L_N}(t)
+\varepsilon_{\mathrm{space}}(N).
\end{equation}
The admissibility and learned-energy statements of Proposition~\ref{prop:structure} remain exact finite-dimensional properties of $\PhiT$; equation \eqref{eq:transfer-bound} does not identify $\cE_\theta$ with the PDE energy $\cE$.
\end{theorem}

\begin{proof}
Set $y(t)=\PhiT(P_Nu_0,t)$ and $x(t)=\widetilde S_t^N(P_Nu_0)$. The two curves have the same initial value. Adding and subtracting $\cFN(y(t))$, using \eqref{eq:matching-error-revised} and the $L_N$-Lipschitz property of $\cFN$, gives
\begin{equation*}
\norm{y(t)-x(t)}_N
\le\int_0^t\bigl(\varepsilon_{\mathrm{match}}+L_N\norm{y(s)-x(s)}_N\bigr)\,ds.
\end{equation*}
Continuous Gronwall yields \eqref{eq:exact-flow-comparison}. Reconstruction stability, the triangle inequality, and \eqref{eq:semidiscrete-conv-revised} then give \eqref{eq:transfer-bound}.
\end{proof}

\begin{remark}[Common-set and stability requirements]\label{rem:transfer-common-set}
The common set $\cV_N$ closes the trajectory-domain requirement in the proof: the matching estimate must hold where the learned trajectory is actually evaluated, not merely at the initial samples.  The exponential factor in \eqref{eq:transfer-bound} uses the Lipschitz constant of the reference semidiscrete generator $\cFN$.  If a one-sided Lipschitz estimate for the learned generator is available, \Cref{thm:generator-stability} gives the sharper three-case factor $\chi_{\omega_N}$, including the uniformly bounded contractive case $\omega_N<0$.
\end{remark}

\begin{proposition}[Adding numerical time-integration errors]\label{prop:numerical-transfer}
Let $t=nh\le T$. On one stated initial-data set, suppose all exact and numerical trajectories remain in the domains of Assumption~\ref{ass:semidiscrete}, $(Q_h^N)^n$ approximates $\widetilde S_t^N$ uniformly for $0\leq nh\leq T$ with error $\eta_N(h,T)$, and a numerical implementation $\widehat\Phi_{\theta,h}^{\,n}$ approximates the exact latent flow uniformly on the same time interval with error $\eta_\theta(h,T)$. Then
\begin{align}\label{eq:numerical-transfer}
\norm{I_N\widehat\Phi_{\theta,h}^{\,n}(P_Nu_0)-\cS_tu_0}_X
&\le C_I\eta_\theta(h,T)
+C_I\varepsilon_{\mathrm{match}}\chi_{L_N}(t)
+\varepsilon_{\mathrm{space}}(N),\\
\norm{I_N(Q_h^N)^n(P_Nu_0)-\cS_tu_0}_X
&\le C_I\eta_N(h,T)+\varepsilon_{\mathrm{space}}(N).
\end{align}
\end{proposition}

\begin{proof}
Insert the corresponding exact finite-dimensional flow in each expression and apply the triangle inequality, Theorem~\ref{thm:transfer}, and reconstruction stability.
\end{proof}

\begin{remark}[Numerical semigroup defect]\label{rem:numerical-semigroup-defect}
For a finite-step implementation, the measured composition residual
\begin{equation*}
\norm{\widehat\Phi_{\theta,h}(\widehat\Phi_{\theta,h}(v,s),t)
-\widehat\Phi_{\theta,h}(v,s+t)}
\end{equation*}
is a \emph{numerical semigroup defect}. It is not the error $\eta_\theta$ itself and is not an exact-flow theorem. The exact autonomous latent flow has zero composition defect, whereas RK4 generally does not. Substep-refinement experiments are therefore needed before attributing a measured residual to integrator accuracy.
\end{remark}

\section{Expressivity within the implemented architecture}\label{sec:expressivity}

The implemented family uses a shared scalar potential, a prescribed interaction graph with Gram-constrained coefficients, and local diagonal mobility. These restrictions preclude an unqualified density claim for arbitrary smooth gradient fields. The following result records what follows rigorously once component approximation has been established for an architecture-compatible target.

\begin{proposition}[Componentwise generator error]\label{prop:universal-approx}
Let $Z\subset\R^N$ be compact and suppose an architecture-compatible target has the form
\begin{equation*}
F^\star(g(z))=-\mathrm{D}g(z)K^\star(z)q^\star(z),
\end{equation*}
where $q^\star=\nabla\Psi^\star$ has the shared-potential and prescribed-interaction form of \eqref{eq:psi-grad}, and $K^\star$ has the permitted local diagonal form. Assume on $Z$ that
\begin{equation*}
\norm{\mathrm{D}g(z)}\le M_g,\quad
\norm{K^\star(z)}\le M_K,\quad
\norm{q^\star(z)}\le M_\Psi,
\end{equation*}
and that one parameter choice satisfies
\begin{equation*}
\sup_Z\norm{K_\theta-K^\star}\le\varepsilon_K,
\qquad
\sup_Z\norm{\nabla\Psi_\theta-q^\star}\le\varepsilon_\Psi.
\end{equation*}
Then
\begin{equation}\label{eq:universal-approx}
\sup_{z\in Z}\norm{F_\theta(g(z))-F^\star(g(z))}
\le M_g\bigl(M_K\varepsilon_\Psi+M_\Psi\varepsilon_K
+\varepsilon_K\varepsilon_\Psi\bigr).
\end{equation}
\end{proposition}

\begin{proof}
Add and subtract $\mathrm{D}g(z)K^\star(z)\nabla\Psi_\theta(z)$ and use submultiplicativity. The bound $\norm{\nabla\Psi_\theta}\le M_\Psi+\varepsilon_\Psi$ gives \eqref{eq:universal-approx}.
\end{proof}

\begin{remark}[Scope of the expressivity result]\label{rem:ua-scope}
Proposition~\ref{prop:universal-approx} is an error-propagation result, not a universal approximation theorem. Simultaneous density under the shared scalar potential, finite-rank Gram interactions, fixed graph, and local diagonal mobility remains open. In particular, the transport term in viscous Burgers is not a pure gradient flow in this model class; Burgers is therefore an out-of-class empirical stress test. A natural extension is a metriplectic or port-Hamiltonian decomposition $\dot z=(J_\theta-K_\theta)\nabla\Psi_\theta$ with $J_\theta^\top=-J_\theta$.
\end{remark}

\begin{corollary}[Conditional architecture-compatible approximation]\label{cor:resolve-prob1}
Consider a sequence of spatial resolutions and model classes for which Assumption~\ref{ass:semidiscrete} holds. Suppose the component errors in \Cref{prop:universal-approx} tend to zero and the factors $M_g$, $M_K$, and $M_\Psi$ remain uniformly bounded on the stated compact sets.  Take the matching error to be the generator bound supplied by \Cref{prop:universal-approx}, namely
\begin{equation*}
\varepsilon_{\mathrm{match}}(N):=M_g(N)\bigl(M_K(N)\varepsilon_\Psi(N)+M_\Psi(N)\varepsilon_K(N)+\varepsilon_K(N)\varepsilon_\Psi(N)\bigr),
\end{equation*}
and assume
\begin{equation*}
\varepsilon_{\mathrm{space}}(N)+C_I(N)\chi_{L_N}(T)\varepsilon_{\mathrm{match}}(N)\longrightarrow0.
\end{equation*}
Then the exact latent flow converges to the PDE semigroup on the stated initial-data set and time interval according to \eqref{eq:transfer-bound}. Any numerical implementation additionally carries the explicit term $\eta_\theta(h,T)$ from \Cref{prop:numerical-transfer}, which must also tend to zero for numerical convergence.
\end{corollary}


\section{Research agenda and open problems}\label{sec:analytical-questions}
The formulation above separates five analytical tasks that are often conflated in one-step operator learning. The first is an approximation question for nonlinear semigroups. The second is an architecture question: which structural properties can be enforced exactly by the model class itself? The third is a boundary-domain question: whether the discrete affine relation is consistent and stable for the intended continuous trace or flux condition. The fourth is a long-time composition question, concerned with how local approximation and numerical composition defects accumulate. The fifth is a transfer question, asking how continuous PDE-level structure is inherited by finite-dimensional semidiscrete and learned generators.

\begin{problem}[Semigroup expressivity]\label{prob:expressivity}
Establish simultaneous component-density results for the exact shared-potential, Gram-constrained interaction, fixed-graph, and local-mobility classes. Proposition~\ref{prop:universal-approx} propagates component errors but does not prove that these errors can be driven to zero for arbitrary gradient-type generators; extending the architecture to non-gradient generators is a further problem.
\end{problem}

\begin{problem}[Exact structural realization]\label{prob:structural-realization}
Characterize neural operator architectures for which admissibility and dissipation are exact model-class properties, rather than penalties enforced only through training data. For such architectures, analogues of \eqref{eq:s4} and \eqref{eq:s5} should hold for all admissible parameters on the stated state set.
\end{problem}

\begin{problem}[Rollout stability under composition defects]\label{prob:rollout-stability}
Develop long-time error estimates for learned evolution families whose semigroup law is only approximate. Given a stability mechanism for repeated time-stepping, quantify how trajectory-level approximation error and semigroup defect accumulate under composition, and determine sharp analogues of \eqref{eq:error-bound} beyond the Lipschitz setting.
\end{problem}

\begin{problem}[Continuous-to-discrete structural transfer]\label{prob:transfer}
For each PDE, verify the theorem card described in \Cref{sec:boundary-semiflow}: continuous well-posedness and the correct temporal object, boundary-domain consistency, method-of-lines convergence, reconstruction stability, a common invariant trajectory set, generator matching, and a stability constant. \Cref{thm:generator-stability,cor:layered-pde-rollout,thm:transfer} provide conditional bridges once these PDE-specific hypotheses are supplied; they do not establish those hypotheses for every benchmark.
\end{problem}

\section{Conclusion}
This paper formulates structure-preserving semigroup learning for nonlinear evolution PDEs and establishes a layered set of conditional results.

\emph{First}, we derive a discrete stability analysis (Proposition~\ref{prop:error-decomposition}) that reveals how one-step approximation error and semigroup defect combine under repeated composition. This analysis (Section~\ref{sec:repeated-composition}) shows that the semigroup defect enters the long-time error bound with a factor of~$n$ (the number of composition steps), making semigroup consistency a first-order concern for long-time rollout, not merely a formal property.

\emph{Second}, we construct a semidiscrete latent dissipative architecture (Sections~\ref{sec:latent-architecture}--\ref{sec:neural-param}) for which invariance of the interior state set $(m,M)^N$, exact semigroup composition, and dissipation of an induced discrete energy are \emph{exact} architectural guarantees---holding for every parameter~$\theta$ and requiring no training penalty---provided the latent ODE generates a global flow (\Cref{prop:global-flow}). These guarantees are verified by construction at the finite-dimensional level and are independent of training quality.

\emph{Third}, for the separate boundary-family architecture, \Cref{thm:boundary-global-semiflow} proves global existence, uniqueness, exact boundary invariance, and the forward semigroup law on each tested discrete affine boundary space. \Cref{prop:explicit-rk-boundary} proves exact-arithmetic stage preservation, while \Cref{thm:rk4-composition-defect} shows that the autonomous RK4 direct-versus-composed discrepancy is a fourth-order numerical quantity under stated hypotheses.

\emph{Fourth}, \Cref{thm:generator-stability} proves that generator matching yields a rollout estimate only after a one-sided stability bound is supplied.  The contractive case prevents indefinite growth of the generator-error contribution.  \Cref{cor:layered-pde-rollout} then separates initial-state, generator, time-integration, reconstruction, and spatial errors.

\emph{Fifth}, the revised Transfer Theorem (\Cref{thm:transfer}) compares the exact latent and exact semidiscrete flows by continuous Gronwall, with explicit projection, reconstruction, norm, common-trajectory, and spatial-convergence hypotheses. \Cref{prop:numerical-transfer} adds numerical time-integration errors without calling them spatial errors or exact semigroup defects. \Cref{prop:universal-approx} then gives a componentwise generator bound for architecture-compatible targets; it is deliberately not a universal approximation theorem.

\emph{Gap between architecture-level and PDE-level guarantees.}
The finite-dimensional guarantees concern either the exact bounded latent flow or the exact boundary-tangent flow; the current implementations do not yet combine every state bound, physical invariant, boundary family, and stability estimate in one architecture.  These results do not imply physical-energy decay or PDE accuracy without the hypotheses of \Cref{thm:generator-stability,thm:transfer}.  Dirichlet, Neumann, and Robin reconstruction is proved for the stated discrete endpoint equations, not as a mesh-convergence theorem for arbitrary continuous boundaries. Component density remains open even within the restricted gradient class (\Cref{prob:expressivity}), while Cahn--Hilliard and transport-dominated systems require conservative or skew extensions. The archived $\beta_V=0$ runs are also outside the coercive global-flow theorem; the revised code exposes a fixed positive floor for new theorem-covered runs.

\emph{Directions for future work.}
Several concrete extensions merit investigation:
\begin{enumerate}[label=(\arabic*)]
    \item \textbf{Two-dimensional spatial domains.} The current architecture is formulated for one-dimensional spatial domains with local nearest-neighbor coupling. Extending to two-dimensional domains requires a spatial interaction structure (e.g., graph neural network coupling or convolutional mobility) that preserves the dissipative gradient-flow structure while scaling to larger state dimensions.
    \item \textbf{Tighter error bounds.} The Transfer Theorem bound scales linearly with $T$ through the factor $e^{L_TT}\,T$. Sharper estimates exploiting the dissipative structure of the latent flow---for example, contractivity of the linearized flow or exponential decay of perturbations---could yield bounds that do not grow (or grow more slowly) with the time horizon.
    \item \textbf{Non-gradient generators.} Adding a skew component, for example $\dot z=(J_\theta-K_\theta)\nabla\Psi_\theta$ with $J_\theta^\top=-J_\theta$, would broaden applicability to PDEs with advection, transport, or non-symmetric dynamics.
    \item \textbf{Adaptive time stepping.} The current formulation uses a fixed step size $\Delta t$. An adaptive scheme that varies $\Delta t$ along the trajectory---exploiting the continuous-time latent flow---could improve efficiency for stiff or multiscale problems.
    \item \textbf{Rigorous convergence rates.} The archived resolution study suggests decreasing error, but an aligned multi-seed rerun and a proof relating spatial resolution, network width, optimization, and PDE regularity are still required.
\end{enumerate}

The framework therefore separates what can be guaranteed exactly at the architecture level from what must be addressed through further analysis, and the research agenda laid out in Section~\ref{sec:analytical-questions} provides a concrete roadmap for closing the gap between finite-dimensional structural guarantees and PDE-level semigroup approximation.

\begin{thebibliography}{9}

\bibitem{ChenRubanovaBettencourtDuvenaud2018NODE}
R.~T.~Q. Chen, Y.~Rubanova, J.~Bettencourt, and D.~K. Duvenaud.
Neural ordinary differential equations.
In \emph{Advances in Neural Information Processing Systems}, 2018.

\bibitem{GreydanusDzambaYosinski2019HNN}
S.~Greydanus, M.~Dzamba, and J.~Yosinski.
Hamiltonian neural networks.
In \emph{Advances in Neural Information Processing Systems}, 2019.

\bibitem{JinZhangZhuTangKarniadakis2020SympNets}
P.~Jin, Z.~Zhang, A.~Zhu, Y.~Tang, and G.~E. Karniadakis.
SympNets: Intrinsic structure-preserving symplectic networks for identifying Hamiltonian systems.
\emph{Neural Networks}, 132:166--179, 2020.

\bibitem{KovachkiLiLiuEtAl2023NeuralOperator}
N.~Kovachki, Z.~Li, B.~Liu, K.~Azizzadenesheli, K.~Bhattacharya, A.~Stuart, and A.~Anandkumar.
Neural operator: Learning maps between function spaces with applications to PDEs.
\emph{Journal of Machine Learning Research}, 24(89):1--97, 2023.

\bibitem{LiKovachkiAzizzadenesheliEtAl2021FNO}
Z.~Li, N.~Kovachki, K.~Azizzadenesheli, B.~Liu, K.~Bhattacharya, A.~Stuart, and A.~Anandkumar.
Fourier neural operator for parametric partial differential equations.
In \emph{International Conference on Learning Representations}, 2021.

\bibitem{LuJinKarniadakis2021DeepONet}
L.~Lu, P.~Jin, and G.~E. Karniadakis.
Learning nonlinear operators via DeepONet based on the universal approximation theorem of operators.
\emph{Nature Machine Intelligence}, 3:218--229, 2021.

\bibitem{YuTianELi2021OnsagerNet}
H.~Yu, X.~Tian, W.~E, and Q.~Li.
OnsagerNet: Learning stable and interpretable dynamics using a generalized Onsager principle.
\emph{Physical Review Fluids}, 6:114402, 2021.

\bibitem{Desjardins2024PHNN}
R.~Desjardins, M.~Pusch, and M.~Biegler.
Port-Hamiltonian neural networks for learning dissipative dynamics.
\emph{arXiv preprint arXiv:2402.03457}, 2024.

\bibitem{Pazy1983Semigroups}
A.~Pazy.
\emph{Semigroups of Linear Operators and Applications to Partial Differential Equations}.
Springer, New York, 1983.

\bibitem{Barbu2010NonlinearDE}
V.~Barbu.
\emph{Nonlinear Differential Equations and Monotone Banach Spaces}.
Springer, Dordrecht, 2010.

\bibitem{CazenaveHaraux1998Semilinear}
T.~Cazenave and A.~Haraux.
\emph{An Introduction to Semilinear Evolution Equations}.
Oxford University Press, Oxford, 1998.

\bibitem{RaissiPerdikarisKarniadakis2019PINNs}
M.~Raissi, P.~Perdikaris, and G.~E. Karniadakis.
Physics-informed neural networks: A deep learning framework for solving forward and inverse problems involving nonlinear partial differential equations.
\emph{Journal of Computational Physics}, 378:686--707, 2019.

\bibitem{E2017Proposal}
W.~E.
A proposal on machine learning via dynamical systems.
\emph{Communications in Mathematics and Statistics}, 5:1--11, 2017.

\bibitem{SirignanoSpiliopoulos2018DGM}
J.~Sirignano and K.~Spiliopoulos.
DGM: A deep learning algorithm for solving partial differential equations.
\emph{Journal of Computational Physics}, 375:1339--1364, 2018.

\bibitem{HundsdorferVerwer2003Numerical}
W.~Hundsdorfer and J.~Verwer.
\emph{Numerical Solution of Time-Dependent Advection-Diffusion-Reaction Equations}.
Springer, Berlin, 2003.

\bibitem{KarlsenRisebro2003Invariant}
K.~H. Karlsen and N.~H. Risebro.
On the uniqueness and stability of entropy solutions of nonlinear degenerate parabolic equations with rough coefficients.
\emph{Discrete and Continuous Dynamical Systems}, 9:1081--1104, 2003.

\bibitem{PareschiRusso2005IMEX}
L.~Pareschi and G.~Russo.
Implicit--explicit numerical methods for stiff hyperbolic systems.
\emph{SIAM Journal on Numerical Analysis}, 43:1291--1310, 2005.

\bibitem{ZhongDeySusmelj2020SympRNN}
Y.~D. Zhong, B.~Dey, and A.~Chakraborty.
Symplectic recurrent neural networks.
In \emph{International Conference on Learning Representations}, 2020.

\bibitem{MassaroliPoliPark2020StableFlows}
S.~Massaroli, M.~Poli, J.~Park, A.~Yamashita, and H.~Asama.
Stable neural flows.
\emph{arXiv preprint arXiv:2003.08563}, 2020.

\bibitem{KidgerMorrilFoster2020NeuralDE}
P.~Kidger, J.~Morrill, J.~Foster, and T.~Lyons.
Neural controlled differential equations for irregular time series.
In \emph{Advances in Neural Information Processing Systems}, 2020.

\bibitem{LiuGomezSuli2022InvariantRegion}
Y.~Liu, J.~G{\'o}mez-Serrano, and A.~Stuart.
Invariant-region-preserving neural networks.
\emph{arXiv preprint arXiv:2206.05441}, 2022.

\bibitem{GeshkovskiRigollet2025NeuralODE}
B.~Geshkovski and P.~Rigollet.
The interplay between discretization and regularization in neural ODEs.
\emph{arXiv preprint arXiv:2501.11663}, 2025.

\bibitem{RuthottoHaber2020DeepPDE}
L.~Ruthotto and E.~Haber.
Deep neural networks motivated by partial differential equations.
\emph{Journal of Mathematical Imaging and Vision}, 62:352--364, 2020.

\bibitem{SanzSerna2016Geometric}
J.~M. Sanz-Serna.
Geometric integration.
In \emph{The Princeton Companion to Applied Mathematics}, pages 153--161. Princeton University Press, 2016.

\bibitem{HairerLubichWanner2006Geometric}
E.~Hairer, C.~Lubich, and G.~Wanner.
\emph{Geometric Numerical Integration: Structure-Preserving Algorithms for Ordinary Differential Equations}.
Springer, Berlin, second edition, 2006.

\bibitem{GottliebShuTadmor2001StrongStability}
S.~Gottlieb, C.-W. Shu, and E.~Tadmor.
Strong stability-preserving high-order time discretization methods.
\emph{SIAM Review}, 43:89--112, 2001.

\bibitem{Hornik1991Approximation}
K.~Hornik.
Approximation capabilities of multilayer feedforward networks.
\emph{Neural Networks}, 4(2):251--257, 1991.

\bibitem{Pinkus1999Approximation}
A.~Pinkus.
Approximation theory of the MLP model in neural networks.
\emph{Acta Numerica}, 8:143--195, 1999.

\bibitem{ChernogorovaTemov2022InvariantRegion}
T.~Chernogorova and R.~Temov.
Invariant-region-preserving schemes for reaction--diffusion systems.
\emph{Preprint}, 2022.

\bibitem{LuLu2020DeepONet}
J.~Lu and L.~Lu.
Universal approximation of solution operators.
\emph{Preprint}, 2020.

\bibitem{Cybenko1989}
G.~Cybenko.
Approximation by superpositions of a sigmoidal function.
\emph{Mathematics of Control, Signals and Systems}, 2(4):303--314, 1989.

\bibitem{CrandallLiggett1971}
M.~G. Crandall and T.~M. Liggett.
Generation of semi-groups of nonlinear transformations on general Banach spaces.
\emph{American Journal of Mathematics}, 93(2):265--298, 1971.

\bibitem{LagarisFotiadis1998}
I.~E. Lagaris, A.~Likas, and D.~I. Fotiadis.
Artificial neural networks for solving ordinary and partial differential equations.
\emph{IEEE Transactions on Neural Networks}, 9(5):987--1000, 1998.

\bibitem{HanJentzenE2018}
J.~Han, A.~Jentzen, and W.~E.
Solving high-dimensional partial differential equations using deep learning.
\emph{Proceedings of the National Academy of Sciences}, 115(34):8505--8510, 2018.

\bibitem{DupontDoucetTeh2019}
E.~Dupont, A.~Doucet, and Y.~W. Teh.
Augmented neural ODEs.
In \emph{Advances in Neural Information Processing Systems}, 2019.

\end{thebibliography}

\end{document}
